\documentclass[11pt,letterpaper]{article}
\usepackage[margin=1in]{geometry}
\usepackage{amsmath,amssymb,amsthm}
\usepackage[T1]{fontenc}
\pdfmapfile{+lm.map}\pdfmapfile{+cm.map}\pdfmapfile{+symbols.map}
\usepackage{lmodern,microtype}
\PassOptionsToPackage{hyphens}{url}
\usepackage[hidelinks]{hyperref}
\usepackage{booktabs,enumitem,longtable,array}
\setlist{itemsep=3pt,topsep=5pt}
\allowdisplaybreaks[2]
\newtheorem{theorem}{Theorem}[section]
\newtheorem{proposition}[theorem]{Proposition}
\newtheorem{lemma}[theorem]{Lemma}
\newtheorem{corollary}[theorem]{Corollary}
\theoremstyle{remark}\newtheorem{remark}[theorem]{Remark}
% Both sources
\newcommand{\E}{\mathbb E}
\newcommand{\dd}{\,\mathrm d}
\newcommand{\ee}{\mathrm e}
% Part I (source 57)
\newcommand{\R}{\mathbb R}
\newcommand{\C}{\mathbb C}
\newcommand{\Res}{\operatorname{Res}}
\newcommand{\Var}{\operatorname{Var}}
% Part II (source 56)
\newcommand{\Log}{\operatorname{Log}}
% Notes added when the two manuscripts were merged into this report
\newenvironment{writenote}{\par\smallskip\begin{quote}\small\textbf{[write, 2 October 2026]}\ }{\end{quote}\smallskip}
\newenvironment{partsource}{\par\medskip\noindent\begin{minipage}{\linewidth}\small\raggedright\textit{Source.}\ }{\end{minipage}\par\medskip}
\newenvironment{partabstract}{\par\begin{quote}\small\textbf{Abstract (as delivered).}\ }{\end{quote}\par}
\expandafter\def\expandafter\UrlBreaks\expandafter{\UrlBreaks\do\_\do-\do/}
\Urlmuskip=0mu plus 1mu
\newcommand{\repopath}[1]{\nolinkurl{#1}}
\hypersetup{pdftitle={Precise asymptotics of the powered Catalan numbers (OEIS A113227)},pdfsubject={Positive Bessel spectrum, all fixed orders, controlled inversion and fixed Fourier sectors for OEIS A113227; merged from two manuscripts of batch 77}}
\title{Precise asymptotics of the powered Catalan numbers (OEIS A113227)\\[4pt]
\large Positive Bessel spectrum, all fixed orders, and fixed Fourier sectors}
\author{A merged research report built from two manuscripts\\
(batch 77, manuscripts 57 and 56)}
\date{Manuscripts dated 1 October 2026; merged 2 October 2026}
\begin{document}
\hypersetup{pageanchor=false}
\maketitle
\begin{abstract}
The powered Catalan numbers $a_n$ count permutations avoiding the vincular pattern $1$-$23$-$4$ (OEIS A113227). This report prints two manuscripts as two Parts. Part~\ref{pcn:part:core} derives a positive discrete moment representation $a_n=\sum_\lambda\omega_\lambda\lambda^n$ from the known formal continued fraction, identifies its transform with a branch-safe Bessel quotient, locates every large spectral node, and proves $a_n\sim\ee^{\ee-2}(2\pi)^{-1/2}r^{-3/2}B_n(\ee)$ with $w=W(n/\ee)$, $r=n/w$ and the Touchard polynomial $B_n$, an expansion to every fixed algebraic order with explicit $c_1$, $c_2$, a canonical continuous inverse and an exponential-generating-function corollary; this answers the question of precise asymptotics that Elizalde left open in 2006. Part~\ref{pcn:part:fs} proves a zero-free complex wedge for the exact Bessel weight, replaces the spectral sum by a lattice sum with relative error $O(\ee^{-dn})$ for every $d<\log2$, and decomposes $a_n$, for every fixed $J$, into $J+1$ exact cutoff-defined Fourier sectors with error bounded by the next complex-saddle envelope; each sector has an all-orders expansion with the same rational coefficients. No convergence of the correction series, growing number of sectors, Borel summation, transseries or exponentially accurate inverse is claimed. The report is unrefereed, and none of it is formalized.
\end{abstract}
\setcounter{tocdepth}{2}
\tableofcontents
\hypersetup{pageanchor=true}

\section*{Guide to this report}
\addcontentsline{toc}{section}{Guide to this report}
\phantomsection\label{pcn:guide}
This guide was written when the two manuscripts were merged. Everything after it, from Part~\ref{pcn:part:core} on, is the manuscripts' text, except for the passages marked \textbf{[write, 2 October 2026]} and the changes listed under ``Conventions of the merge'' below.

\subsection*{What the report proves}
\phantomsection\label{pcn:guide:results}
\begin{itemize}
\item \textbf{Part~\ref{pcn:part:core}} (source 57). Lemma~\ref{pcn:lem:finite} and Section~\ref{pcn:sec:moments}: every finite truncation of the continued fraction is a probability moment transform, and the limit measure $\mu$ is unique, with $a_n=\int s^n\dd\mu(s)$. Theorem~\ref{pcn:thm:spectrum}: the transform is $1-E/D$ with entire $E$, $D$, all zeros of $D$ are positive and simple, and they carry positive weights $\omega_\lambda$. Proposition~\ref{pcn:prop:nodes}: one node $\lambda_k=k+O(\ee^{-ck\log k})$ near every large integer, with an exact weight formula. Theorem~\ref{pcn:thm:main}: the equivalent \eqref{pcn:eq:leading} and the expansion \eqref{pcn:eq:allorders} to every fixed order, with $c_1$, $c_2$ and the algorithm \eqref{pcn:eq:c-algorithm}. Theorem~\ref{pcn:thm:inverse} and \eqref{pcn:eq:inversefirst}--\eqref{pcn:eq:inversesecond}: the canonical continuous inverse. Corollary~\ref{pcn:cor:EGF}: the exponential generating function is entire, with an all-orders expansion on the positive axis.
\item \textbf{Part~\ref{pcn:part:fs}} (source 56). Lemma~\ref{pcn:fs:lem:wedge}: the exact Bessel weight $Q(z)=1/(\pi^2zb(z)^2)$ has an asymptotic expansion, and $b$ is eventually zero-free, in every closed wedge $|\arg z|\le\theta<\pi/2$. Proposition~\ref{pcn:fs:prop:replace}: global spectral replacement with relative error $O_d(\ee^{-dn})$, $d<\log 2$. Theorem~\ref{pcn:fs:thm:main}: the fixed-sector decomposition \eqref{pcn:fs:eq:decomp}, the all-orders expansion \eqref{pcn:fs:eq:sector-expansion} of every fixed exact sector at the complex saddle $w_j=W_{-j}(n/\ee)$, and the envelope ordering \eqref{pcn:fs:eq:envelope}.
\end{itemize}

\subsection*{Sources and provenance}
\phantomsection\label{pcn:guide:provenance}
The report was built from two manuscripts, both dated 1 October 2026, both delivered in the arrival commit \texttt{096ee7b87} of 2 October 2026 as batch~77, manuscripts 57 and 56, and placed in this directory by commit \texttt{aa7345800} (batch 77P2). Neither pins a ProveIt commit or names a repository path. Their author lines are ``A proof and reproducible research report'' (source 57) and ``A research addendum with reproducible diagnostics'' (source 56); no tool or person is named.

\begin{center}
\footnotesize\setlength{\tabcolsep}{4pt}
\begin{tabular}{@{}llp{0.55\linewidth}lll@{}}
\toprule
Part & Source & Archive and main file & Pages & Labels & Prefix\\
\midrule
I & 57 & \texttt{oeis-powered-catalan-report.zip},\newline \texttt{powered-catalan-asymptotics.tex} (635 lines) & 14 & 59 & \texttt{pcn:}\\
II & 56 & \texttt{oeis-powered-catalan-fourier-addendum.zip},\newline \texttt{powered-catalan-fourier-addendum.tex} (497 lines) & 10 & 32 & \texttt{pcn:fs:}\\
\bottomrule
\end{tabular}
\end{center}

\noindent\textbf{Nesting.} Source 56 is an addendum of source 57 and is not self-contained: it cites source 57 (bibliography key \texttt{Core}) for the spectrum, the real coefficient algorithm and the interpolation, and its archive embeds source 57's manuscript, PDF and verification record byte for byte (as \texttt{core-report.tex}, \texttt{core-report.pdf}, \texttt{core-mathematical-verification.md}), together with six byte-identical support files (\texttt{residue\_expansion.py}, \texttt{saddle\_expansion.py}, their three outputs and \texttt{requirements.txt}). The placement staged each file once, from source 57's archive. Source 56 does not repeat any proof of source 57.

\noindent\textbf{Where the merge had to choose.} (i) The Parts follow the dependency order 57, 56. (ii) Both sources are printed in full; nothing of either is omitted except the title blocks and the two bibliographies, which are merged below. (iii) Source 56 restates the inputs it takes from source 57 (the weight and node estimates \eqref{pcn:fs:eq:spectral-input}--\eqref{pcn:fs:eq:weight-bound}, the coefficients $c_0$, $c_1$, $c_2$, $\alpha_0$, $\alpha_1$, $\alpha_2$, the generators \eqref{pcn:fs:eq:beta-generator}, \eqref{pcn:fs:eq:alpha-generator} and \eqref{pcn:fs:eq:coefficient-generator}). These restatements are kept, because Part~\ref{pcn:part:fs} continues them into the complex wedge and its proofs cite them by its own labels; each now names the place in Part~\ref{pcn:part:core} where the input is proved. (iv) Three symbols of source 56 were renamed to avoid collisions with Part~\ref{pcn:part:core}; see the notation table.

\subsection*{How the Parts depend on one another}
\phantomsection\label{pcn:guide:dependencies}
Part~\ref{pcn:part:fs} imports from Part~\ref{pcn:part:core} the positive discrete spectrum (Theorem~\ref{pcn:thm:spectrum}), the node and weight estimates (Proposition~\ref{pcn:prop:nodes}), the real coefficient algorithm \eqref{pcn:eq:c-algorithm} and the interpolation \eqref{pcn:eq:interpolation}. It proves everything else itself. The sector expansion \eqref{pcn:fs:eq:sector-expansion} at $j=0$ is the all-orders expansion \eqref{pcn:eq:allorders} of Part~\ref{pcn:part:core} for the exact principal sector $I_0(n)$, whereas Part~\ref{pcn:part:core} expands $a_n$ itself. The two statements differ by the nonzero sectors, which are smaller than every algebraic order (Theorem~\ref{pcn:fs:thm:main}), so they are consistent. Part~\ref{pcn:part:fs} does not change the inverse of Part~\ref{pcn:part:core} and proves no exponentially accurate inverse.

\subsection*{Relation to other reports}
\phantomsection\label{pcn:guide:relations}
\begin{writenote}
\emph{Sibling report on the pattern $12$-$34$.} Batch 77 also opened the report \repopath{SetTheory/Cardinals/docs/reports/generating-functions-and-asymptotics/oeis-sequence-asymptotics/a113226-vincular-avoiders}, on permutations avoiding the vincular pattern $12$-$34$ (OEIS A113226). Both reports prove precise asymptotics for a length-four vincular pattern whose asymptotics Elizalde studied in 2006 \cite{Elizalde}, and both cite Callan's bijection paper on $1$-$23$-$4$ \cite{Callan}. Otherwise they share no theorem and no method: the A113226 report works from a closed exponential generating function (an inverse-square-root singularity and a stretched-exponential saddle in powers of $n^{-1/3}$), this report from a positive Bessel spectrum and a Touchard saddle. The two sequences are different, as Section~\ref{pcn:sec:intro} says. They were therefore kept as two reports, by the coordinator's decision at the write, rather than printing these two manuscripts as further Parts of the A113226 report.
\end{writenote}
The Touchard saddle $w=W(n/\ee)$ of this report is the analogue, at parameter $\theta=\ee$, of the Lambert saddle $r=W_0(n)$ of the Bell numbers ($\theta=1$). The Bell numbers are treated in the neighbouring report \repopath{a277364-bell-asymptotics} of this collection directory and in the chapter ``The Bell numbers: a Lambert saddle'' of the transseries volume \repopath{Analysis/Transseries/docs/series-and-transseries/Transseries_And_Inversion/} (label \texttt{q2:sec:bell}; all-orders theorem \texttt{q2:thm:bell}). Remark \texttt{q2:rem:bell-sectors} of that volume keeps the nonprincipal complex saddles of the Bell numbers \emph{formal}. Part~\ref{pcn:part:fs} proves fixed complex sectors rigorously, but for A113227, a different sequence with an exact Bessel weight; it does not upgrade that remark (see the note at the end of Part~\ref{pcn:part:fs}).

\subsection*{Placement}
\phantomsection\label{pcn:guide:placement}
Part~\ref{pcn:part:fs} is the batch's closest case to the repository's rule that transseries material goes to \repopath{Analysis/Transseries/}. It was kept in this collection, beside its foundation, because it proves a finite decomposition into a fixed number of exact cutoff-defined Fourier sectors with algebraic expansions and constructs no transseries, Borel sum or Stokes data; its own scope paragraph (Section~\ref{pcn:fs:sec:scope}) says so. If Vladimir prefers every exponentially ordered sector theorem in the transseries tree, Part~\ref{pcn:part:fs} can be filed there whole and replaced here by a cross-reference.

\subsection*{Notation}
\phantomsection\label{pcn:guide:notation}
Each Part keeps its source's letters, except for the three renamed symbols marked below. Symbols not in the table are local to the Part that uses them. The table prints the tempting false reading beside the true one.

{\small
\begin{longtable}{@{}>{\raggedright\arraybackslash}p{0.15\linewidth}>{\raggedright\arraybackslash}p{0.46\linewidth}>{\raggedright\arraybackslash}p{0.33\linewidth}@{}}
\toprule
Symbol & Meaning and where & Do not read it as\\
\midrule
\endfirsthead
\toprule
Symbol & Meaning and where & Do not read it as\\
\midrule
\endhead
\bottomrule
\endfoot
$a_n$ & The count (both Parts); in Part~\ref{pcn:part:fs}, for real $n$, the interpolation $\mathcal A(n)$ of \eqref{pcn:eq:interpolation}. & Part~\ref{pcn:part:core}'s Bessel combination $a(t)$ \eqref{pcn:eq:ab}; Part~\ref{pcn:part:fs}'s window endpoint $a=\lfloor r/8\rfloor+\frac12$.\\
$w$, $r$ & $w=W(n/\ee)$ (real branch), $r=n/w=\ee^{w+1}$ (both Parts). & \\
$w_j$, $t_j$ & Part~\ref{pcn:part:fs}: $w_j=W_{-j}(n/\ee)$, $t_j=n/w_j$, the complex saddles; $w_0=w$, $t_0=r$. & Part~\ref{pcn:part:core}'s $w_0$, $r_0$, $N_0$ in \eqref{pcn:eq:coreinverse} solve the inverse core equation $\ee^{w_0+1}(w_0^2+1)=L$; that $w_0$ is \emph{not} $W(n/\ee)$.\\
$M_n$, $M_j$, $E_j$ & $M_n$ of \eqref{pcn:eq:rw} (Part~\ref{pcn:part:core}) is $M_0$ of Part~\ref{pcn:part:fs}; $M_j$ is the complex leading factor at $t_j$ and $E_j=|M_j|$ the envelope. & $E_j$ is not Part~\ref{pcn:part:core}'s entire function $E(t)$ \eqref{pcn:eq:E}, nor the exponential generating function $\mathcal E(z)$.\\
$E_{\mathrm S}(z)$ & Part~\ref{pcn:part:fs}: the remainder in Schl\"afli's representation \eqref{pcn:fs:eq:schlafli}. \emph{Renamed} from source 56's $E(z)$. & Not $E(t)$ of \eqref{pcn:eq:E}.\\
$B_n(\theta)$ & Touchard (Bell) polynomial (Part~\ref{pcn:part:core}); $B_n(\ee)$ is not the Bell number. & Part~\ref{pcn:part:fs}'s upper window endpoint $B=\lfloor4r\rfloor+\frac12$.\\
$\beta_{2m}$ & Bernoulli numbers (both Parts). Part~\ref{pcn:part:fs}'s generator \eqref{pcn:fs:eq:alpha-generator} writes $\beta_{2q}$; \emph{renamed} from source 56's $B_{2q}$. & Not a coefficient of $\mathcal B$.\\
$\mathcal B(z)$, $\mathcal B_m$ & The formal series \eqref{pcn:eq:Bformal} $=$ \eqref{pcn:fs:eq:beta-generator}, and its coefficients $\mathcal B_m=[s^m]\mathcal B(s)$: $1,-\frac12,-\frac5{24},-\frac7{16},\ldots$, see \eqref{pcn:eq:Bcoeff}. Part~\ref{pcn:part:fs}'s $\mathcal B_m$ is \emph{renamed} from source 56's $\beta_m$. & $\mathcal B_2=-\frac5{24}$ is not the Bernoulli number $\beta_2=\frac16$.\\
$\alpha_\ell$ & Coefficients of the weight expansion, \eqref{pcn:eq:alpha} $=$ \eqref{pcn:fs:eq:alpha-generator} (both Parts; Part~\ref{pcn:part:fs} indexes them by $m$). & \\
$c_j$ & The rational correction functions (Part~\ref{pcn:part:core}: $c_j(w)$; Part~\ref{pcn:part:fs}: $c_\ell(u)$, the same functions). & The constant $0<c<1$ of \eqref{pcn:eq:nodeerror} and \eqref{pcn:fs:eq:spectral-input}.\\
$q(t)$, $Q(z)$ & Part~\ref{pcn:part:core}: $q=a/b$, a small ratio. Part~\ref{pcn:part:fs}: $Q(z)=1/(\pi^2zb(z)^2)$, the exact weight, so $\omega_k=Q(\lambda_k)(\ldots)$, cf.\ \eqref{pcn:eq:weight-b}. & $Q$ is not a multiple of $q$.\\
$J$, $L$ & Part~\ref{pcn:part:core}: $J$ is the algebraic order in \eqref{pcn:eq:allorders} and of $H_J$; $L$ is a log-target in Section~\ref{pcn:sec:inverse}. Part~\ref{pcn:part:fs}: $J$ is the number of retained nonzero Fourier sectors, $L$ the algebraic order of each sector. & Part~\ref{pcn:part:fs}'s $J$ is not an algebraic order.\\
$H$, $h$ & Part~\ref{pcn:part:core}: $H(t)$, the Bessel denominator \eqref{pcn:eq:H}; $H_J(x)$, the inverse \eqref{pcn:eq:HJ}. Part~\ref{pcn:part:fs}: $H=\lfloor R^{1/4}\rfloor$, a cutoff; $H_j$, horizontal contour integrals; $\mathcal H(c)=c-1-\log c$; $h(x)=x-1-\log x$; $h=\Im t_m$, a contour height. & \\
$I_j(n)$, $\mathcal I(z)$ & Part~\ref{pcn:part:fs}: exact Fourier sectors \eqref{pcn:fs:eq:sectors}; the rotated Schl\"afli integral \eqref{pcn:fs:eq:rotated}. & Not the modified Bessel function $I_\nu$ of Part~\ref{pcn:part:core}, Section~\ref{pcn:sec:bessel}.\\
$\sigma$ & Part~\ref{pcn:part:core}: $\sqrt{r/(w+1)}$. Part~\ref{pcn:part:fs}: $\sqrt{r/w}$. & The two widths differ.\\
$d$, $K$, $R$ & Part~\ref{pcn:part:core}: horizontal weights $d_h$, the inverse constant $d_0$; a Poisson variable $K$; $R_k(t)$, and $R_n=a_n/M_n$. Part~\ref{pcn:part:fs}: a decay rate $d<\log2$; a large fixed integer $K$; $R=|z|$. & \\
$P$, $F$, $\theta$ & Part~\ref{pcn:part:core}: $P(t)=t^{t/2}/\Gamma(t+1)$, the transform $F(t)$, $F_0(x)$, the Touchard parameter $\theta$. Part~\ref{pcn:part:fs}: reciprocal products $P_{j,\delta}$, the tilted exponent $F_p$, a wedge angle $\theta$. & \\
\end{longtable}
}

\subsection*{Status, scope and non-claims}
\phantomsection\label{pcn:guide:status}
Neither manuscript states how it was produced; both reached the repository through its incoming channel for AI-assisted research deliveries and are treated as AI-assisted. Both are unrefereed: source 57's verification record calls itself ``not a proof-assistant verification or an external peer-review certification'', and source 56 calls its review ``component cross-review plus integrated mathematical review, not formal verification or conventional peer review''. Their computations are high-precision diagnostics, not interval certificates. Nothing in this report is formalized in Lean or Rocq, and placing it in this repository confers no formal status.

The union of the sources' non-claims, all kept in the Parts, is as follows.
\begin{itemize}
\item No convergence of the infinite correction series (Theorem~\ref{pcn:thm:main}, Part~\ref{pcn:part:fs} Section~\ref{pcn:fs:sec:scope}); the all-orders result is a Poincar\'e expansion, not an exponentially complete transseries.
\item Fixed truncation only: the number of sectors and the algebraic order are fixed; no growing-$J$ theorem, no summation of infinitely many sector expansions, no canonical Borel sum. A fixed algebraic truncation of the principal sector has an error larger than every nonzero Fourier envelope (Remark~\ref{pcn:fs:rem:scope}).
\item The error scales are proved without explicit constants. A numerical Newton residual controls only the root of the truncated inverse equation; a certified enclosure of the true inverse needs an explicit remainder bound. No exponentially accurate inverse is claimed. A finite expansion in inverse logarithms is not claimed to give integer accuracy uniformly.
\item The numerical tables are diagnostics. Part~\ref{pcn:part:fs}'s quadratures integrate central saddle segments and do not evaluate complete cutoff-defined sectors.
\item The literature checks are scoped and do not establish worldwide priority or repository-wide nonduplication. The recurrence, path models and continued fraction are prior art (Callan, Beaton et al., Gladkovskii in OEIS); Poisson summation, complex Gaussian expansion and Lambert-$W$ saddles in Touchard asymptotics are standard (Paris).
\end{itemize}

\subsection*{Instances of results elsewhere in the repository}
\phantomsection\label{pcn:guide:lambert}
Several steps are instances of results proved in the transseries volume \repopath{Analysis/Transseries/docs/series-and-transseries/Transseries_And_Inversion/}. No novelty is claimed for these mechanics.
\begin{itemize}
\item The saddle $w=W(n/\ee)$ solves $w+\log w=\log(n/\ee)$, the linear--logarithmic equation $aX+b\log X=L$ with $a=b=1$; the volume's Lambert core, theorem \texttt{p0:thm:lambert-core}, solves it on every branch. Part~\ref{pcn:part:fs}'s complex saddles $w_j=W_{-j}(n/\ee)$, with $w_j+\Log w_j=\log(n/\ee)-2\pi ij$, are the branches $k=-j$ of the same solution.
\item Theorem~\ref{pcn:thm:inverse} and the explicit inverses \eqref{pcn:eq:inversefirst}--\eqref{pcn:eq:inversesecond} revert a perturbed core equation, as in theorem \texttt{p0:thm:perturbed-inversion}. The threshold rule $\min\{n:a_n\ge Y\}=\lceil x_*(\log Y)\rceil$, with a comparison of neighbouring exact values when an enclosure straddles an integer, is the separation step of theorem \texttt{p0:thm:staircase}, whose separation inequality is formalized in Lean as \texttt{staircase\_separation} in \repopath{Analysis/FabiusFunction/Lean/FabiusFunction/StaircaseInversion.lean}. Nothing specific to A113227 is formalized.
\end{itemize}

\subsection*{Conventions of the merge}
\phantomsection\label{pcn:guide:conventions}
\begin{itemize}
\item Every label carries the report prefix: \texttt{pcn:} in Part~\ref{pcn:part:core} and \texttt{pcn:fs:} in Part~\ref{pcn:part:fs}, for example \texttt{pcn:fs:thm:main}. Guide labels are \texttt{pcn:guide:}, Part labels \texttt{pcn:part:}, and the few section labels added at the merge \texttt{pcn:sec:} and \texttt{pcn:fs:sec:}.
\item Sections are numbered through the report: source 57's Section $k$ is Section $k$ here (Sections 1--10), source 56's is Section $k+10$ (Sections 11--17). Theorems are numbered within sections, as in both sources, so source 56's Theorem 1.1 is Theorem 11.1. Equations are numbered through the report.
\item Source 56's four citations of source 57 (\texttt{\textbackslash cite\{Core\}}) and its two further mentions of it (``the audited spectrum'', ``the companion report'') are now references to Part~\ref{pcn:part:core}, and the bibliography entry \texttt{Core}, which described \texttt{core-report.pdf} and \texttt{core-report.tex} as included in the package, is dropped. The two \texttt{DLMF} entries are merged into one.
\item Sentences about ``the reproducibility package'', ``the reproducibility archive'', ``the accompanying output'' or ``the supplied symbolic generator'' describe the delivered archives, which survive in commit \texttt{096ee7b87}. Which of their files this report ships, and under which names, is listed in the README. The delivered PDFs, checksum manifests and source 56's embedded copies are not shipped.
\item Each Part opens with its source's abstract as delivered.
\end{itemize}

\cleardoublepage
\part{Precise asymptotics of the powered Catalan numbers}\label{pcn:part:core}

\begin{partsource}
Batch-77 manuscript 57, archive \texttt{oeis-powered-catalan-report.zip}, main file \texttt{powered-catalan-asymptotics.tex} (635 lines, 14-page PDF), dated 1 October 2026. Author line: ``A proof and reproducible research report''. No pin. Printed in full, with labels prefixed \texttt{pcn:}; four section labels and the \textbf{[write]} notes were added at the merge.
\end{partsource}

\begin{partabstract}
The powered Catalan numbers count permutations avoiding the vincular pattern $1$-$23$-$4$ and form OEIS A113227. We derive a positive discrete moment representation from their known formal continued fraction, identify its transform by a branch-safe Bessel quotient, and determine every sufficiently large spectral node and its weight. This proves a sharp equivalent involving the Bell polynomial $B_n(\ee)$, including its explicit constant, and gives an expansion to every fixed algebraic order. An exact coefficient algorithm, two displayed corrections, and a canonical continuous inverse accompany the proof. The analytic arguments control the full spectrum and all omitted tails; numerical computations are used only as reproducible diagnostics.
\end{partabstract}

\section{Main results and known starting points}\label{pcn:sec:intro}
Let $a_n$ count permutations $\pi_1\cdots\pi_n$ for which there are no indices $i<j<j+1<k$ with
\[
 \pi_i<\pi_j<\pi_{j+1}<\pi_k.
\]
Thus $a_n$ is OEIS A113227, with
\[
 (a_n)_{n\ge0}=1,1,2,6,23,105,549,3207,20577,143239,\ldots.
\]
The adjacency of the middle pair is essential. This sequence is different from avoidance of the classical pattern $1234$, fully consecutive $1234$, or the vincular pattern $12$-$34$.

For $n>0$, throughout the report put
\begin{equation}\label{pcn:eq:rw}
 w=W(n/\ee),\qquad r=\frac n w=\ee^{w+1},\qquad
 M_n=\frac{\ee^{-2}r^{n-3/2}\exp(r-n)}{\sqrt{2\pi(w+1)}},
\end{equation}
where $W$ is the positive real Lambert function. Write
\[
 B_n(\theta)=\sum_{k=0}^n\left\{\!\begin{matrix}n\\k\end{matrix}\!\right\}\theta^k
\]
for the Bell or Touchard polynomial; $B_n(\ee)$ is not the ordinary Bell number.

\begin{theorem}\label{pcn:thm:main}
The powered Catalan numbers satisfy
\begin{equation}\label{pcn:eq:leading}
 a_n\sim M_n\sim
 \frac{\ee^{\ee-2}}{\sqrt{2\pi}}\,r^{-3/2}B_n(\ee).
\end{equation}
For every fixed integer $J\ge1$,
\begin{equation}\label{pcn:eq:allorders}
 a_n=M_n\left(\sum_{j=0}^{J-1}\frac{c_j(w)}{r^j}
                    +O_J(r^{-J})\right),
\end{equation}
where $c_0=1$ and the $c_j$ are explicitly computable rational functions. Their denominators divide a constant times $(w+1)^{3j}$, and their numerators have degree at most $3j$. In particular,
\begin{align}
 c_1(w)&=\frac{20w^3+86w^2+128w+67}{24(w+1)^3},\label{pcn:eq:c1}\\
 c_2(w)&=\frac{1168w^6+9968w^5+34692w^4+64536w^3
                   +69380w^2+42208w+11857}{1152(w+1)^6}.\label{pcn:eq:c2}
\end{align}
The constants in \eqref{pcn:eq:allorders} may depend on $J$; no convergence of the infinite correction series is asserted.
\end{theorem}

The recurrence, combinatorial models, and formal fraction used below are known. Callan \cite{Callan} gives the recurrence and valley-marked Dyck paths; Beaton et al.\ \cite{Beaton} develop the powered-Catalan terminology and further models. OEIS credits the formal continued fraction to Sergei N.\ Gladkovskii, 10 October 2012 \cite{OEIS}. Elizalde \cite{Elizalde} established asymptotic bounds and explicitly left a precise asymptotic for this pattern open. The contribution proved here is the sharp equivalent, the global positive spectral representation supporting it, and the controlled refinements. A scoped literature review is summarized in Section~\ref{pcn:sec:literature}; it does not certify worldwide priority or repository-wide nonduplication.
\begin{writenote}
Part~\ref{pcn:part:fs} (source 56) refines Theorem~\ref{pcn:thm:main} beyond all algebraic orders. For every fixed $J$ it writes $a_n$ as the sum of $J+1$ exact cutoff-defined Fourier sectors, with an error bounded by the next complex-saddle envelope $E_{J+1}$, and each fixed sector has an all-orders expansion with the same rational functions $c_j$, evaluated at the complex saddle $w_j=W_{-j}(n/\ee)$ (Theorem~\ref{pcn:fs:thm:main}). The nonzero sectors are of relative size $\exp(-2\pi^2j^2r/(w+1)+\cdots)$, so they do not affect \eqref{pcn:eq:allorders}.
\end{writenote}

\section{The formal fraction and positive finite measures}
\subsection{A reversible path derivation}
In the known path model, a Dyck path of semilength $n$ has $h+1$ possible marks at every valley of height $h$. Keep mark-zero valleys unchanged. Replace a valley $DU$ with a nonzero mark by a long horizontal step at height $h+1$, with $h$ possible colors. Distinct valleys do not overlap, and expansion of each horizontal step to $DU$ reverses the operation. The resulting Schr\"oder paths have unit up/down weights and horizontal weights
\[
 d_0=0,\qquad d_h=h-1\quad(h\ge1).
\]
First-return decomposition therefore gives, as a formal series,
\begin{equation}\label{pcn:eq:cf}
 A(x)=\sum_{n\ge0}a_nx^n=1+\frac{x}{U_0(x)-x},\qquad
 U_k(x)=1-kx-\frac{x}{U_{k+1}(x)}\quad(k\ge0).
\end{equation}
Equivalently it is the Thron fraction with successive denominator weights $0,0,1,2,\ldots$ and unit numerators. This also falls under general nonnegative Thron-fraction moment-positivity results \cite{PSZ}; the following direct proof is sufficient here. Nothing in the argument treats the divergent series $A(x)$ as analytic near zero.

\subsection{Finite fractions}
Truncate \eqref{pcn:eq:cf} at $U_m=1-mx$, with $m\ge1$, and denote the resulting rational function by $A_m$. For $k\ge1$ let $R_k(t)=tU_k(1/t)$ and at the root let $R_0=tU_0(1/t)-1$. Then
\begin{equation}\label{pcn:eq:R}
 R_m=t-m,\quad R_k=t-k-\frac{t}{R_{k+1}}\ (1\le k<m),
 \quad R_0=t-1-\frac{t}{R_1},\quad A_m(1/t)=1+\frac1{R_0}.
\end{equation}
The exceptional root shift must be retained.

\begin{lemma}\label{pcn:lem:finite}
Each $A_m(1/t)$ is a probability moment transform
\[
 A_m(1/t)=\sum_{j=1}^{N_m}\omega_{m,j}\frac{t}{t-\lambda_{m,j}},
 \qquad \omega_{m,j}>0,\quad\lambda_{m,j}>0,\quad
 \sum_j\omega_{m,j}=1.
\]
\end{lemma}
\begin{proof}
Consider a real rational function
\[
 R(t)=t-b-\sum_j\frac{c_j}{t-\eta_j},\qquad c_j>0,\quad\eta_j>0,
 \qquad R(0)<0.
\]
Repeated poles can be combined. Since $R'(t)=1+\sum_j c_j/(t-\eta_j)^2>0$, its zeros are simple and interlace its poles. There is no nonpositive zero, one zero between zero and its first pole, one between each successive pair, and one after its last pole. Its reciprocal consequently has the partial fraction expansion
\[
 \frac1{R(t)}=\sum_j\frac{\rho_j}{t-\xi_j},\qquad
 \rho_j=\frac1{R'(\xi_j)}>0,\qquad \sum_j\rho_j=1.
\]
The sum of residues follows from $1/R(t)\sim1/t$ at infinity. Hence
\[
 \frac{t}{R(t)}=1+\sum_j\frac{\xi_j\rho_j}{t-\xi_j}.
\]
This proves preservation of the stated form under \eqref{pcn:eq:R}, starting from $t-m$. Also $R_k(0)=-k$ for $k\ge1$ and $R_0(0)=-1$. At the root evaluation at zero gives $\sum_j\rho_j/\xi_j=1$, and
\[
 1+\frac1{R_0(t)}=\sum_j\frac{\rho_j}{\xi_j}\frac{t}{t-\xi_j}.
\]
These are the asserted probability weights.
\end{proof}

\subsection{Passage to a unique measure}\label{pcn:sec:moments}
Every fixed formal coefficient of $A_m$ eventually stabilizes to that of $A$. The associated probability measures $\mu_m$ have bounded first moments, so they are tight. Along any weakly convergent subsequence, a bound on the next moment gives uniform integrability of each fixed power. The limit $\mu$ thus has moments
\begin{equation}\label{pcn:eq:moments}
 a_n=\int_0^\infty s^n\dd\mu(s),\qquad n\ge0.
\end{equation}
It is unique: the permutation interpretation gives $a_n\le n!$, so monotone convergence implies
\[
 \int\exp(\delta s)\dd\mu(s)
   =\sum_{n\ge0}\frac{\delta^na_n}{n!}\le\frac1{1-\delta}
       \quad(0<\delta<1).
\]
The moments determine the moment-generating function near zero and therefore determine the measure. Thus the whole sequence $\mu_m$ converges to $\mu$. In particular the function
\begin{equation}\label{pcn:eq:Fmeasure}
 F(t)=\int_0^\infty\frac{t}{t-s}\dd\mu(s),\qquad
 t\in\C\setminus[0,\infty),
\end{equation}
is the limit of the finite fractions. It is this analytic transform, rather than the divergent ordinary series, that will be identified.

\section{A branch safe Bessel representation}\label{pcn:sec:bessel}
\subsection{Identification on the negative real axis}
Fix $t=-s<0$ and define
\[
 V_k=\frac{I_{k+s-1}(2\sqrt{s})}{\sqrt{s}\,I_{k+s}(2\sqrt{s})}.
\]
All these quantities are positive: denominator orders are positive and numerator orders exceed $-1$. The modified Bessel recurrence gives
\begin{equation}\label{pcn:eq:Irec}
 V_k=1+\frac{k}{s}+\frac{1/s}{V_{k+1}},\qquad V_k\ge1+\frac{k}{s}.
\end{equation}
The truncated $U_k(-1/s)$ satisfy the same recurrence, with terminal value $1+m/s$. Its discrepancy from $V_m$ is at most $(1/s)/(1+(m+1)/s)$. Each preceding step contracts the discrepancy by at most
\[
 \frac{1/s}{(1+(k+1)/s)^2}=\frac{s}{(s+k+1)^2}\le\frac1{4(k+1)}.
\]
Their product tends to zero factorially. This proves the required branch selection directly, without an assumption about a minimal solution of an infinite recurrence.

In Bessel notation the resulting transform is
\begin{equation}\label{pcn:eq:BesselF}
 F(t)=1-\frac{J_{-t}(2\sqrt t)}{\sqrt t\,J_{-t-1}(2\sqrt t)+J_{-t}(2\sqrt t)}.
\end{equation}
A formula with no square-root or logarithmic branch convention is preferable for global use. Define
\begin{align}
 E(t)&=\sum_{j\ge0}\frac{(-t)^j}{j!\,\Gamma(j-t+1)},\label{pcn:eq:E}\\
 D(t)&=\sum_{j\ge0}\frac{(j-t+1)(-t)^j}{j!\,\Gamma(j-t+1)}.\label{pcn:eq:D}
\end{align}
The gamma growth in $j$ proves local uniform convergence for complex $t$, so both functions are entire. The Bessel series \cite[\S10.2]{DLMF} identifies \eqref{pcn:eq:BesselF}, after cancellation of its common factor $t^{-t/2}$, with
\begin{equation}\label{pcn:eq:entirequotient}
 F(t)=1-\frac{E(t)}{D(t)}.
\end{equation}
The negative-axis equality and the identity theorem give its analytic continuation.

\begin{theorem}\label{pcn:thm:spectrum}
All zeros of $D$ are positive and simple. They are exactly the support points of a probability measure with positive weights, and
\begin{equation}\label{pcn:eq:spectrum}
 F(t)=\sum_{\lambda:D(\lambda)=0}\omega_\lambda\frac{t}{t-\lambda},
 \qquad
 a_n=\sum_{\lambda:D(\lambda)=0}\omega_\lambda\lambda^n,
 \qquad
 \omega_\lambda=-\frac{E(\lambda)}{\lambda D'(\lambda)}>0.
\end{equation}
The support has no finite accumulation point and is bounded away from zero.
\end{theorem}
\begin{proof}
The functions $E,D$ have no common zero. At zero both equal 1. At a nonzero common zero, \eqref{pcn:eq:BesselF} would make two consecutive Bessel functions vanish at the same nonzero argument. Their recurrence would make every forward consecutive order vanish, contradicting the fixed-argument large-positive-order series.

The quotient \eqref{pcn:eq:entirequotient} agrees with the Stieltjes transform \eqref{pcn:eq:Fmeasure}; hence its genuine poles can occur only on the positive axis. Every such pole is simple, since a Cauchy transform of a finite positive measure has at most a simple meromorphic pole, and its residue is the corresponding positive atom. Between poles the quotient continues analytically with real values. Stieltjes inversion therefore gives zero measure there, ruling out continuous mass. Entire-function zeros are locally finite. Finally $F(0)=0$ rules out an atom at zero, while $D(0)=1$ leaves a neighborhood of zero free of support. This proves every assertion, including the absence of an additional entire contribution or a hidden cut.
\end{proof}

\section{Uniform estimates and the large spectral nodes}
For real $t>0$ define
\begin{align}
 a(t)&=J_t(2\sqrt t)-\sqrt t\,J_{t+1}(2\sqrt t),\label{pcn:eq:ab}\\
 b(t)&=\sqrt t\,Y_{t+1}(2\sqrt t)-Y_t(2\sqrt t),\qquad q(t)=\frac{a(t)}{b(t)}.
\end{align}
Order reflection \cite[\S10.4]{DLMF} writes the denominator in \eqref{pcn:eq:BesselF} as
\begin{equation}\label{pcn:eq:H}
 H(t)=\cos(\pi t)a(t)+\sin(\pi t)b(t).
\end{equation}

\begin{lemma}\label{pcn:lem:largeorder}
Let $P(t)=t^{t/2}/\Gamma(t+1)$. As real $t\to\infty$,
\begin{align}
 J_t(2\sqrt t)&=P(t)\bigl(\ee^{-1}+O(t^{-1})\bigr),\label{pcn:eq:Jest}\\
 b(t)&=-\frac{\ee}{\pi}\,t\Gamma(t)t^{-t/2}\bigl(1+O(t^{-1})\bigr),\label{pcn:eq:best}\\
 a(t)&=O(P(t)),\quad a'(t)=O(P(t)\log t),\quad
 b'(t)=O(|b(t)|\log t).\label{pcn:eq:derivs}
\end{align}
Consequently, for every fixed $0<c<1$ and all sufficiently large $t$,
\begin{equation}\label{pcn:eq:qsmall}
 |q(t)|+|q'(t)|\le\exp(-ct\log t).
\end{equation}
\end{lemma}
\begin{proof}
The absolutely convergent series gives
\[
 \frac{J_t(2\sqrt t)}{P(t)}
 =\sum_{j\ge0}\frac{(-1)^j}{j!}
                 \prod_{h=1}^j(1+h/t)^{-1}.
\]
The product differs from 1 by at most $j(j+1)/(2t)$, whose sum against $1/j!$ is finite. Its derivative in $t$ is bounded by $j(j+1)/(2t^2)$. Since $P'/P=O(\log t)$, the same reasoning for the adjacent order proves the $J$ and $a$ estimates without differentiating an unspecified remainder.

Schl\"afli's formula \cite[10.9.7]{DLMF} is
\begin{align*}
 Y_\nu(z)={}&\frac1\pi\int_0^\pi\sin(z\sin\theta-\nu\theta)\dd\theta\\
 &-\frac1\pi\int_0^\infty
       \bigl(\ee^{\nu u}+\cos(\pi\nu)\ee^{-\nu u}\bigr)
                                      \ee^{-z\sinh u}\dd u.
\end{align*}
The bounded integral and the decaying integral contribute $O(\sqrt t)$ to $b$ and to its total $t$-derivative. In the growing integral use $v=\sqrt t\,\ee^u$. This gives the exact dominant integral
\begin{equation}\label{pcn:eq:bint}
 b(t)=-\frac{t^{-t/2}}\pi
 \int_{\sqrt t}^\infty(v-1)v^{t-1}\exp(-v+t/v)\dd v+O(\sqrt t).
\end{equation}
For $V$ of gamma distribution with shape $t$ and scale 1, its integral normalized by $\Gamma(t)$ is
\[
 \E\bigl[(V-1)\ee^{t/V}\mathbf1_{\{V\ge\sqrt t\}}\bigr]
       =\ee t\bigl(1+O(t^{-1})\bigr).
\]
Indeed on $[t/2,2t]$ Taylor expansion at $V=t$ uses $\E(V-t)=0$ and $\Var(V)=t$. On the lower tail $\ee^{t/V}\le\ee^{\sqrt t}$, absorbed by the gamma bound $\ee^{-c_0t}$; on the upper tail it is at most $\ee^{1/2}$. This proves \eqref{pcn:eq:best}.

Differentiation of the growing integrand including its prefactor adds
\[
 \log v-\frac12\log t-\frac12+\frac1v.
\]
The same gamma localization bounds its weighted integral by $O(\log t)$ times the main integral. The lower-endpoint term is $O(t^{-1/2})$ in absolute size and is negligible relative to that main term. This proves the last derivative bound. Stirling's formula now gives \eqref{pcn:eq:qsmall}.
\end{proof}

\begin{proposition}\label{pcn:prop:nodes}
For every sufficiently large integer $k$ there is exactly one spectral node $\lambda_k$ in $(k-\tfrac12,k+\tfrac12)$, there are no other sufficiently large nodes, and
\begin{equation}\label{pcn:eq:nodeerror}
 \lambda_k=k+O\bigl(\exp(-ck\log k)\bigr)\qquad(0<c<1).
\end{equation}
At those nodes the weight has the exact expression
\begin{equation}\label{pcn:eq:weight-exact}
 \omega_k=\frac1{\lambda_kb(\lambda_k)^2
     \{\pi^2(1+q(\lambda_k)^2)+\pi q'(\lambda_k)\}}.
\end{equation}
In particular,
\begin{align}
 \omega_k&=\frac1{\pi^2k b(k)^2}
       \bigl(1+O(\exp(-ck\log k))\bigr),\label{pcn:eq:weight-b}\\
 \omega_k&\sim\ee^{-2}\frac{k^{k-3}}{\Gamma(k)^2}
     \sim\frac{k^{-2}\exp(-k\log k+2k)}{2\pi\ee^2}
     \sim\frac{\ee^{-2}}{\sqrt{2\pi}}k^{-3/2}\frac{\ee^k}{k!}.
 \label{pcn:eq:weight-leading}
\end{align}
\end{proposition}
\begin{proof}
For large $t$, $b(t)<0$. Zeros of \eqref{pcn:eq:H} solve $\tan(\pi t)+q(t)=0$. On each indicated interval the left side ranges from $-\infty$ to $+\infty$ and has derivative at least $\pi-|q'|>0$. The endpoints are not zeros because $b\ne0$ there. This proves completeness and \eqref{pcn:eq:nodeerror}.

At a pole, reflection gives the residue formula
\[
 \omega_k=-\frac{J_t+q(t)Y_t}
 {t b(t)\{\pi(1+q(t)^2)+q'(t)\}}\bigg|_{t=\lambda_k},
\]
where the argument of each Bessel function is $2\sqrt t$. The Wronskian \cite[10.5.2]{DLMF} implies
\[
 J_tb+aY_t=\sqrt t\,(J_tY_{t+1}-J_{t+1}Y_t)=-1/\pi.
\]
This proves \eqref{pcn:eq:weight-exact}. Bounds \eqref{pcn:eq:derivs} and \eqref{pcn:eq:qsmall} permit replacing $\lambda_k$ by $k$, giving \eqref{pcn:eq:weight-b}. Equations \eqref{pcn:eq:best} and Stirling's formula finish the proof.
\end{proof}

A further expansion, not needed for completeness of the spectrum, gives
\[
 a(t)=P(t)\ee^{-1}\left(t^{-2}+\frac32t^{-3}-\frac{97}{24}t^{-4}
                               +O(t^{-5})\right).
\]
It follows that the displacement is eventually positive and
\begin{equation}\label{pcn:eq:shift}
 \lambda_k-k\sim\ee^{-2}\frac{k^{k-4}}{\Gamma(k)^2}\sim\frac{\omega_k}{k}.
\end{equation}
The finite number of remaining low nodes cause no difficulty in large moments.

\section{The residue expansion to arbitrary order}
Introduce the following \emph{formal} series:
\begin{align}
 T_0(z)&=\sum_{j\ge0}\frac1{j!}\prod_{h=0}^{j-1}(1-hz)^{-1},\label{pcn:eq:T0}\\
 T_1(z)&=\sum_{j\ge0}\frac1{j!}\prod_{h=1}^{j}(1-hz)^{-1},\qquad
 \mathcal B(z)=\ee^{-1}\{T_0(z)-zT_1(z)\}.\label{pcn:eq:Bformal}
\end{align}
Each formal coefficient is a convergent Poisson sum of a polynomial in $j$. The expressions are not infinite convergent sums to be evaluated at $z=1/k$: individual factors then have poles for large $j$.

\begin{lemma}\label{pcn:lem:bexp}
For each fixed order, the following is an asymptotic expansion with a remainder of the corresponding inverse power:
\begin{equation}\label{pcn:eq:bexp}
 b(k)\sim-\frac{\ee k\Gamma(k)k^{-k/2}}\pi\,\mathcal B(1/k).
\end{equation}
Its first coefficients are
\begin{align}
 \mathcal B(z)&=1-\frac z2-\frac{5z^2}{24}-\frac{7z^3}{16}
                    -\frac{7057z^4}{5760}+O(z^5),\label{pcn:eq:Bcoeff}\\
 \mathcal B(z)^{-2}&=1+z+\frac76z^2+2z^3+\frac{1739}{360}z^4+O(z^5).
\end{align}
\end{lemma}
\begin{proof}
Use \eqref{pcn:eq:bint} with $t=k$. Fix $0<\epsilon<1/2$ and retain terms $0\le j\le m=\lfloor k^\epsilon\rfloor$ in the exponential $\exp(k/v)$. On $[k/2,2k]$ the discarded remainder is at most a constant times $2^{m+1}/(m+1)!$, smaller than every inverse power of $k$. The exterior contribution of the original integral is exponentially small as in Lemma~\ref{pcn:lem:largeorder}.

For retained $j$ the gamma integrals have shapes $k-j+1$ and $k-j$. Extending their central intervals to $(0,\infty)$ changes them by uniformly exponentially small relative errors, since their centers differ from $k$ by $O(k^\epsilon)$. The extended integral of the $j$th term is
\[
 \frac{k^j}{j!}\bigl\{\Gamma(k-j+1)-\Gamma(k-j)\bigr\}.
\]
After normalization by $k\Gamma(k)$, these terms give precisely $T_0(1/k)-k^{-1}T_1(1/k)$, truncated at $j=m$.

For either reciprocal product, expansion through degree $J$ has remainder bounded uniformly for $j\le m$ by
\[
 C_J k^{-J-1}(1+j)^{2J+2}\exp(Cj^2/k).
\]
This follows by expanding its logarithm or by differentiating the product. The exponential factor is bounded for $\epsilon<1/2$, and the polynomial factor is summable against $1/j!$. Extending every coefficient sum to all $j$ costs a further superalgebraic factorial tail. Thus the remainder is $O_J(k^{-J-1})$ after retaining degrees through $J$, proving \eqref{pcn:eq:bexp}. The coefficients follow by exact polynomial arithmetic.
\end{proof}

Let $\beta_{2m}$ denote the Bernoulli numbers, to distinguish them from $B_n(\theta)$. Define $\alpha_\ell$ by
\begin{equation}\label{pcn:eq:alpha}
 \sum_{\ell\ge0}\alpha_\ell z^\ell
 =\mathcal B(z)^{-2}
   \exp\left(-\sum_{m\ge1}\frac{\beta_{2m}z^{2m-1}}{m(2m-1)}\right).
\end{equation}
Combining \eqref{pcn:eq:weight-b}, Lemma~\ref{pcn:lem:bexp}, and Stirling's formula proves
\begin{equation}\label{pcn:eq:weight-expansion}
 \omega_k=\frac{k^{-2}\exp(-k\log k+2k)}{2\pi\ee^2}
 \left(\sum_{\ell=0}^{J-1}\frac{\alpha_\ell}{k^\ell}+O_J(k^{-J})\right).
\end{equation}
In particular
\[
 \alpha_0=1,\quad \alpha_1=\frac56,\quad \alpha_2=\frac{73}{72},\quad
 \alpha_3=\frac{11821}{6480},\quad \alpha_4=\frac{702533}{155520}.
\]
For comparison with Bell polynomials the equivalent normalization is
\begin{equation}\label{pcn:eq:Dobinskiweight}
 \frac{\omega_k}{[\ee^{-2}/\sqrt{2\pi}]k^{-3/2}\ee^k/k!}
 \sim1+\frac{11}{12k}+\frac{313}{288k^2}
                         +\frac{98959}{51840k^3}+\cdots.
\end{equation}
An exact implementation is included in the reproducibility package. Polynomial moments are evaluated using $\ee^{-1}\sum_{j\ge0}j^p/j!=B_p(1)$, so these coefficients are not numerical fits.

\section{Global moment reduction and the saddle expansion}
\subsection{Removing the node displacements}
The exact positive moment representation separates into finitely many low nodes and the $\lambda_k$ of Proposition~\ref{pcn:prop:nodes}. For $k\ge\sqrt r$, \eqref{pcn:eq:nodeerror} gives, for every fixed $J$,
\[
 \left(\frac{\lambda_k}{k}\right)^n=1+o(r^{-J})
\]
uniformly in this range. For nodes indexed by $k<\sqrt r$, including all low nodes, positivity and total mass one bound their entire contribution by $(2\sqrt r)^n$ for large $n$. This is smaller than every algebraic multiple of a single term near $k=r$, since that term has logarithm $n\log r-n+r+O(\log r)$. Thus the actual nodes can be replaced by integers at every algebraic order.

Consequently \eqref{pcn:eq:weight-expansion} reduces the problem to real lattice sums with exponent
\begin{equation}\label{pcn:eq:phi}
 \phi(t)=n\log t-t\log t+2t.
\end{equation}
The unique maximum is at $t=r$, where
\[
 \phi(r)=n\log r-n+r,\qquad -\phi''(r)=\frac{w+1}{r}.
\]
The Gaussian width $\sigma=\sqrt{r/(w+1)}$ tends to infinity.

\subsection{Uniform localization and remainders}
The exact identity
\begin{equation}\label{pcn:eq:localization}
 \phi(ry)-\phi(r)
 =-r\{w(y-1-\log y)+(y\log y-y+1)\}
\end{equation}
controls both tails: both bracketed functions are nonnegative and vanish only at $y=1$. It also absorbs every fixed power coming from amplitudes and their derivatives.

With $t=r+\sigma u$, the exponent becomes
\begin{align}
 \phi(t)-\phi(r)&=-u^2/2+
          \sum_{h\ge1}\gamma_h(w)r^{-h/2}u^{h+2},\label{pcn:eq:gaussianexp}\\
 \gamma_h(w)&=
 \frac{(-1)^{h+1}((h+1)w+1)}{(h+1)(h+2)(w+1)^{(h+2)/2}}.
\end{align}
For each fixed $h$, these coefficients are bounded for $w\ge1$, as are the normalized amplitude coefficients. On the symmetric interval $|u|\le r^\delta$, with a sufficiently small fixed $\delta>0$, Taylor expansion through degree $2J-1$ in $r^{-1/2}$ has, after including the Gaussian factor, a remainder bounded by
\[
 C_Jr^{-J}P_J(|u|)\ee^{-u^2/4},
\]
where $P_J$ is a fixed polynomial. Its integral is $O_J(r^{-J})$. Odd terms integrate to zero. The rest of the fixed relative neighborhood is Gaussian-small, and \eqref{pcn:eq:localization} handles the exterior. This proves a uniform algebraic expansion without residual powers of $\log r$.

To pass from integrals to sums, apply Euler--Maclaurin to
\[
 f_\ell(t)=t^{-\ell-2}\ee^{\phi(t)}.
\]
For fixed $m,\ell$, its endpoint derivatives vanish once $n$ is large enough, and Gaussian rescaling and the derivative product rule give
\begin{equation}\label{pcn:eq:EMderivative}
 \int_0^\infty|f_\ell^{(2m)}(t)|\dd t
 \le C_{m,\ell}\sigma^{-2m}\ee^{\phi(r)}r^{-\ell-2}\sigma.
\end{equation}
On the exterior, \eqref{pcn:eq:localization} absorbs the finitely many derivative powers of $n$ and $1/t$; near zero take $n>\ell+2+2m$. Hence the relative lattice error is $O_m(((w+1)/r)^m)$. Taking $m=J+1$ makes it $o(r^{-J})$, since $w=\log r-1$. There is no fractional-power lattice correction.

The leading Gaussian integral is
\[
 \frac{1}{2\pi\ee^2}r^{-2}\ee^{\phi(r)}\sqrt{\frac{2\pi r}{w+1}}=M_n.
\]
The preceding bounds applied to the residual term in \eqref{pcn:eq:weight-expansion} prove \eqref{pcn:eq:allorders}.

\subsection{An explicit algorithm for every coefficient}
Let $V$ be a centered Gaussian variable of variance $1/(w+1)$, used only as a device for exact polynomial moments. Then the coefficients in Theorem~\ref{pcn:thm:main} are
\begin{align}\label{pcn:eq:c-algorithm}
 c_j(w)=[z^{2j}]\E\Bigg[&\sum_{\ell\ge0}\alpha_\ell z^{2\ell}
                                      (1+Vz)^{-\ell-2}\notag\\[-2pt]
 &\times\exp\left\{\sum_{h\ge1}
 \frac{(-1)^{h+1}((h+1)w+1)}{(h+1)(h+2)}V^{h+2}z^h\right\}\Bigg].
\end{align}
For fixed $j$ all sums truncate finitely. Use $\E V^{2m}=(2m-1)!!/(w+1)^m$ and $\E V^{2m+1}=0$. Counting the Gaussian powers gives the denominator and degree bounds asserted in the theorem.

For example the amplitude $t^{-2}$ and the cubic/quartic exponent terms give
\[
 c_1=\frac56+\frac3{w+1}-\frac{2w+1}{(w+1)^2}
 -\frac{3w+1}{4(w+1)^2}+\frac{5(2w+1)^2}{24(w+1)^3},
\]
which simplifies to \eqref{pcn:eq:c1}; the next coefficient is \eqref{pcn:eq:c2}. The package gives $c_3,c_4$ and the exact generator.

\subsection{The Bell polynomial form}
Dobi\'nski's identity reads
\[
 B_n(\ee)=\ee^{-\ee}\sum_{k\ge0}\frac{k^n\ee^k}{k!}.
\]
Stirling's formula and the same saddle give
\[
 B_n(\ee)\sim\frac{\ee^{-\ee}r^n\exp(r-n)}{\sqrt{w+1}}.
\]
Comparing with $M_n$ proves the second equivalent in \eqref{pcn:eq:leading}, with constant
\[
 \frac{\ee^{\ee-2}}{\sqrt{2\pi}}=0.818193265211777\ldots.
\]
This completes the proof of Theorem~\ref{pcn:thm:main}.

\section{Canonical continuous inversion}\label{pcn:sec:inverse}
\subsection{The increasing branch}
The spectral measure gives a canonical interpolation
\begin{equation}\label{pcn:eq:interpolation}
 \mathcal A(z)=\sum_{\lambda:D(\lambda)=0}\omega_\lambda\lambda^z,
 \qquad \lambda^z=\exp(z\log\lambda).
\end{equation}
It is entire in $z$: the support has a positive minimum and all positive moments, so the series is locally uniformly dominated on each vertical strip. For real $x$, $\mathcal A(x)>0$, and
\[
 (\log\mathcal A)''(x)=\Var_x(\log\lambda)>0,
\]
where the variance is under the tilted probability weights. The measure is nondegenerate since $a_1=1$ and $a_2=2$. Jensen's inequalities give
\[
 \mathcal A'(0)=\E\log\lambda<0,\qquad
 \mathcal A'(1)=\E(\lambda\log\lambda)>0.
\]
Thus the unique minimum lies in $(0,1)$; $\mathcal A$ is strictly increasing on $[1,\infty)$, from $\mathcal A(1)=1$ to infinity. Let $x_*(L)\ge1$ be its inverse defined by $\mathcal A(x_*)=\ee^L$ for $L\ge0$. All preceding asymptotic arguments hold for real $x\to\infty$.

\subsection{Controlled inversion of the full expansion}
Write $r=r(x)$ and $w=w(x)$ as in \eqref{pcn:eq:rw}, and set
\begin{align}
 F_0(x)&=x\log r-x+r=\ee^{w+1}(w^2+1),\label{pcn:eq:F0}\\
 g(x)&=\frac32\log r+\frac12\log(w+1)+2+\frac12\log(2\pi).
\end{align}
Let $\ell_j(w)=[z^j]\log(\sum_{k\ge0}c_k(w)z^k)$, and define
\begin{equation}\label{pcn:eq:HJ}
 H_J(x)=F_0(x)-g(x)+\sum_{j=1}^J\frac{\ell_j(w(x))}{r(x)^j}\quad(J\ge0).
\end{equation}
Then $\log\mathcal A(x)=H_J(x)+O_J(r^{-J-1})$ and
\[
 H_J'(x)=w+1+O\bigl(1/[r(w+1)]\bigr)>0
\]
for large $x$.

\begin{theorem}\label{pcn:thm:inverse}
If $x_J(L)$ solves $H_J(x_J)=L$ on its eventual increasing branch, then
\begin{equation}\label{pcn:eq:inverse-error}
 x_*(L)=x_J(L)+O_J\left(\frac{r(x_J)^{-J-1}}{w(x_J)+1}\right)
       \qquad(L\to\infty).
\end{equation}
\end{theorem}
\begin{proof}
No differentiation of an asymptotic remainder is needed. Convexity bounds $(\log\mathcal A)'(x)$ between the secants on $[x-1,x]$ and $[x,x+1]$. The already proved expansion on those intervals gives
\[
 (\log\mathcal A)'(x)=w+1+O(1/r).
\]
At $x_J$ the logarithmic residual is $O_J(r^{-J-1})$. The mean-value theorem and the eventual positive lower bound on the derivative prove \eqref{pcn:eq:inverse-error}.
\end{proof}

\subsection{An explicit inverse based on the exact leading core}
Let $w_0>0$ solve
\begin{equation}\label{pcn:eq:coreinverse}
 \ee^{w_0+1}(w_0^2+1)=L,\qquad r_0=\ee^{w_0+1},\qquad N_0=r_0w_0.
\end{equation}
This core equation is strictly increasing for $w_0>0$. Define
\begin{equation}\label{pcn:eq:d0}
 d_0=\frac32+
 \frac{\tfrac12\log(w_0+1)+2+\tfrac12\log(2\pi)}{w_0+1}.
\end{equation}
Since $F_0'(x)=w+1$, Taylor expansion yields
\begin{equation}\label{pcn:eq:inversefirst}
 x_*(L)=N_0+d_0+O\bigl(1/[r_0(w_0+1)]\bigr).
\end{equation}
A further explicit term is
\begin{align}\label{pcn:eq:inversesecond}
 x_*(L)={}&N_0+d_0-\frac{1}{r_0(w_0+1)}
 \left[c_1(w_0)+\frac{d_0^2}{2(w_0+1)}\right.\notag\\[-2pt]
 &\left.\hspace{28mm}-d_0\left\{\frac{3}{2(w_0+1)}
                              +\frac{1}{2(w_0+1)^2}\right\}\right]
 +O\bigl(1/[r_0^2(w_0+1)]\bigr).
\end{align}
Indeed $F_0''=1/[r(w+1)]$ and
$g'=r^{-1}\{3/[2(w+1)]+1/[2(w+1)^2]\}$; substituting $N_0+d_0$ leaves exactly the displayed logarithmic residual of order $r_0^{-1}$. All coefficients and the needed derivatives remain bounded for $w_0\ge1$.

For numerical work, $2W(\tfrac12\sqrt{L/\ee})$ initializes safeguarded Newton for \eqref{pcn:eq:coreinverse}; its derivative is $\ee^{w+1}(w+1)^2$. Newton can then solve \eqref{pcn:eq:HJ}. A numerical residual $\eta$ bounds the numerical error in the $H_J$ root by order $|\eta|/(w+1)$. A certified enclosure for the \emph{true} inverse at a finite target additionally requires an explicit numerical bound for the asymptotic remainder. The $O_J$ estimates above prove error scales, not ready-made interval certificates.

For $Y>1$, the discrete threshold is $\min\{n:a_n\ge Y\}=\lceil x_*(\log Y)\rceil$. If a finite-target enclosure straddles an integer, neighboring exact recurrence values must be compared. A finite expansion purely in inverse logarithms is not claimed to give integer accuracy uniformly.

\section{The exponential generating function on the positive axis}
Let $\mathcal E(z)=\sum_{n\ge0}a_nz^n/n!$, distinguished from the interpolation $\mathcal A$ and the normalized Bessel numerator $E$.

\begin{corollary}\label{pcn:cor:EGF}
The function $\mathcal E$ is entire and has the exact representation
\[
 \mathcal E(z)=\sum_\lambda\omega_\lambda\exp(\lambda z).
\]
For real $z\to+\infty$, put $q=\exp(z+1)$. Then
\begin{equation}\label{pcn:eq:EGF}
 \mathcal E(z)=\frac{\ee^{-2}}{\sqrt{2\pi}}\ee^q q^{-3/2}
 \left(1+\frac{67}{24q}+\frac{11857}{1152q^2}+O(q^{-3})\right).
\end{equation}
There is an expansion to every fixed inverse power of $q$; its coefficients are $c_j(0)$ from \eqref{pcn:eq:c-algorithm}.
\end{corollary}
\begin{proof}
The weight estimate \eqref{pcn:eq:weight-leading} makes the exponential spectral sum locally uniformly convergent for complex $z$. Its Taylor coefficients are the exact moments, proving entirety and the representation. For $k\ge\sqrt q$, the pole displacement in $\exp(z\lambda_k)$ is smaller than every algebraic relative error. The complementary spectrum contributes at most $\exp(2z\sqrt q)$ for large $q$, negligible compared with $\ee^q q^{-3/2}$.

Using \eqref{pcn:eq:Dobinskiweight}, the main sum is
\[
 \frac{\ee^{-2}}{\sqrt{2\pi}}\ee^q
 \E\left[K^{-3/2}\left(1+\frac{11}{12K}
                          +\frac{313}{288K^2}+\cdots\right)
                    \mathbf1_{\{K\ge1\}}\right],
\]
where $K$ is Poisson with mean $q$. Poisson concentration and Taylor expansion about $q$ justify this to every fixed order. For general fixed $p>0$,
\[
 \E[K^{-p}\mathbf1_{\{K\ge1\}}]
 =q^{-p}\left(1+\frac{p(p+1)}{2q}
 +\frac{p(p+1)(p+2)(3p+5)}{24q^2}+O(q^{-3})\right).
\]
The first correction is $15/8+11/12=67/24$; the next is
$665/128+(11/12)(35/8)+313/288=11857/1152$.
Alternatively the real saddle of $(z+2)t-t\log t$ is $q$. Retain the lower cutoff $t\ge\sqrt q$ (or a smooth equivalent): the uncut inverse-power amplitude is not integrable at zero when $w=0$. All cutoff endpoint terms are $\exp(O(\sqrt q\log q))$, negligible relative to $\ee^q$. The central Gaussian coefficient and Taylor bounds remain bounded at $w=0$, and give exactly $c_j(0)$ from \eqref{pcn:eq:c-algorithm}. Euler--Maclaurin on the cut-off domain has negligible endpoints and proves the stated all-orders form.
\end{proof}

\section{Numerical checks and reproducibility}\label{pcn:sec:checks}
The exact recurrence used for validation is
\begin{equation}\label{pcn:eq:recurrence}
 p_{0,0}=1,\qquad p_{n,0}=0\ (n>0),\qquad
 p_{n,k}=p_{n-1,k-1}+k\sum_{j\ge k}p_{n-1,j},\qquad
 a_n=\sum_kp_{n,k}.
\end{equation}
Suffix sums evaluate a row in linear time, using exact integers. Let $R_n=a_n/M_n$. The following high-precision diagnostics use exact $a_n$ through $n=2500$:
\begin{center}
\begin{tabular}{r r r r}
\toprule
$n$ & $R_n$ & $r(R_n-1)$ & $c_1(w)$\\
\midrule
50   &1.058106438359&1.352791674642&1.251473654579\\
100  &1.032686382461&1.240030857333&1.186048117956\\
200  &1.018353693325&1.165057005282&1.135943310080\\
400  &1.010254878707&1.112627221355&1.096857105828\\
800  &1.005694089255&1.074346527794&1.065802240138\\
1600 &1.003140908882&1.045321581963&1.040701211579\\
2500 &1.002134095265&1.030095275457&1.026990044234\\
\bottomrule
\end{tabular}
\end{center}
The second scaled residual $r^2(R_n-1-c_1/r)$ likewise approaches $c_2(w)$; its discrepancy decreases from about $0.244$ at $n=50$ to $0.00638$ at $n=2500$.

A separate spectral computation with 110 decimal digits and 65 nodes begins with
\begin{center}
\begin{tabular}{r r r}
\toprule
 & node $\lambda$ & probability weight $\omega$\\
\midrule
0&0.320955976594004859&0.541638005518968697\\
1&1.115304216978757830&0.256514749784783314\\
2&2.042059059595086477&0.118696820250535227\\
3&3.014766264330034697&0.051578277597176440\\
\bottomrule
\end{tabular}
\end{center}
The sum of these 65 weights agrees with 1 at the displayed 50 digits in the accompanying output. Their moments agree with the exact recurrence with relative errors from about $1.46\times10^{-65}$ at $n=1$ to $1.84\times10^{-40}$ at $n=30$, consistent with the omitted positive tail. This root search is a numerical diagnostic; the analytic proof, rather than the scan, establishes completeness of the spectrum.

The inverse check solves $H_J(x)=\log a_n$ for $J=0,1,2$ and exact targets through $n=500$. At $n=500$, the $J=2$ approximation differs from 500 by $2.07456\times10^{-7}$. Its error multiplied by $r^3(w+1)$ is $2.18743$, compared with the next logarithmic coefficient $2.13558$.

The reproducibility archive contains the editable \LaTeX\ source, a portable build script, exact coefficient generators, all three numerical tests, their outputs, dependency versions, a single replay command, and mathematical verification records. It contains no third-party paper copies. Numerical arithmetic uses \texttt{mpmath}; symbolic coefficients use \texttt{SymPy}. The formulas are derived symbolically, never fitted to these data.
\begin{writenote}
In this repository the five programs and their outputs are shipped under \texttt{code/57-catalan-*} and \texttt{data/57-catalan-*}, with the verification and validation records as \texttt{57-catalan-*.md}; the delivered PDF and checksum manifest are not shipped, and the delivered \LaTeX\ source is the text of this Part. The README lists every file and gives a rerun recipe that leaves the shipped outputs untouched.
\end{writenote}

\section{Literature scope and further questions}\label{pcn:sec:literature}
Elizalde's Proposition 6.1 and Corollary 6.2 \cite{Elizalde} give coarse bounds and subfactorial growth; Section 7 leaves precise asymptotics for $1$-$23$-$4$ open. Callan \cite{Callan}, Duchi--Guerrini--Rinaldi \cite{Duchi}, and Beaton et al.\ \cite{Beaton} supply the models used here. Later work by Lin--Fu \cite{LinFu}, Cervetti \cite{Cervetti}, Frosini--Guerrini--Rinaldi \cite{Frosini}, and Testart \cite{Testart} provides refinements, generating trees, or related bijections. The sources surveyed do not supply a previous precise equivalent of the form \eqref{pcn:eq:leading}. As of the dated review, the OEIS entry lists the known fraction but no sharp asymptotic. These observations are a scoped literature statement, not a proof of global novelty.
\begin{writenote}
In this repository, the neighbouring report \repopath{a113226-vincular-avoiders} (in the same collection directory) proves precise asymptotics for the vincular pattern $12$-$34$ (OEIS A113226), a different sequence, by a closed exponential generating function; it also cites Elizalde's bounds and Callan's $1$-$23$-$4$ bijection paper. The two reports share no theorem or method and are kept separate; see the Guide, ``Relation to other reports''. The Bell-number analogue of the saddle used here is treated in \repopath{a277364-bell-asymptotics} and in the repository's transseries volume (Guide, ``Relation to other reports'').
\end{writenote}

The formal fraction, general continued-fraction techniques, Bessel identities, and saddle methods are established tools. The central analytic issue addressed by this report is their connection to the \emph{correct global positive measure}; a formal quotient alone would not justify the moment asymptotics.

Several further questions remain meaningful.
\begin{enumerate}
\item Can the exact Bessel nodes and weights be enclosed by efficient certified interval algorithms with explicit tail bounds? This would turn the asymptotic inverse estimates into finite-target enclosures.
\item What is the optimally truncated remainder of the algebraic expansion, and how do exponentially small node displacements interact with lattice effects? The present result is an all-orders Poincar\'e expansion, not an exponentially complete transseries.
\item Can the positive spectral representation be deformed to refined powered-Catalan statistics, with a uniform two-parameter asymptotic theorem?
\item Is there a direct probabilistic or combinatorial explanation for the Bell-polynomial parameter $\ee$ and the factor $r^{-3/2}$?
\end{enumerate}
\begin{writenote}
Question~2 is now partly answered. Part~\ref{pcn:part:fs} treats the lattice effects at every fixed exponential order: after the node displacements are removed with relative error $O(\ee^{-dn})$, $d<\log2$ (Proposition~\ref{pcn:fs:prop:replace}), Poisson summation turns the lattice sum into exact Fourier sectors, and any fixed number of them, each with its own all-orders expansion, approximates $a_n$ up to the next envelope (Theorem~\ref{pcn:fs:thm:main}). Still open: the optimally truncated remainder, any statement with a growing number of sectors, a Borel sum, and an exponentially complete transseries, all of which Part~\ref{pcn:part:fs} disclaims (Section~\ref{pcn:fs:sec:scope}). Questions~1, 3 and~4 are untouched; Part~\ref{pcn:part:fs}'s numerical checks are not certified either.
\end{writenote}


\part{Fixed Fourier sectors for the powered Catalan numbers}\label{pcn:part:fs}

\begin{partsource}
Batch-77 manuscript 56, archive \texttt{oeis-powered-catalan-fourier-addendum.zip}, main file \texttt{powered-catalan-fourier-addendum.tex} (497 lines, 10-page PDF), dated 1 October 2026. Author line: ``A research addendum with reproducible diagnostics''. No pin. An addendum to source 57 (Part~\ref{pcn:part:core}), printed in full, with labels prefixed \texttt{pcn:fs:}; its citations of source 57 are now references to Part~\ref{pcn:part:core}, and three symbols are renamed (Guide, ``Notation''): $\mathcal B_m$ for its $\beta_m$, $\beta_{2q}$ for its Bernoulli numbers $B_{2q}$, and $E_{\mathrm S}$ for its Schl\"afli remainder $E$.
\end{partsource}

\begin{partabstract}
The positive Bessel spectrum of the powered Catalan numbers permits a fixed-mode Fourier refinement beyond their all-orders algebraic expansion. We prove a zero-free complex wedge for the exact Bessel weight, replace the full spectral sum by a lattice sum with relative error $O(\ee^{-dn})$ for every $d<\log2$, and control the entire omitted Fourier tail by the next complex-saddle envelope. Each retained exact sector has its own all-orders expansion with the same rational coefficients as the real saddle. Exact sectors and finite asymptotic truncations are kept distinct: a fixed algebraic approximation to the principal sector does not have exponentially improved accuracy.
\end{partabstract}

\section{Definitions and the fixed sector theorem}
Let $a_n$ be OEIS A113227, beginning $1,1,2,6,23,105,549,\ldots$. We use the positive discrete spectral representation established in Part~\ref{pcn:part:core} (Theorem~\ref{pcn:thm:spectrum} and Proposition~\ref{pcn:prop:nodes}). Set
\[
 w=W_0(n/\ee),\qquad r=n/w=\ee^{w+1},
\]
and, with principal logarithms and square roots in the right half-plane, define
\[
 b(z)=\sqrt z\,Y_{z+1}(2\sqrt z)-Y_z(2\sqrt z),\qquad
 Q(z)=\frac1{\pi^2z b(z)^2}.
\]
The letter $Q$ distinguishes the exact weight from Lambert $W$. Use the explicit half-integer endpoints
\[
 a=\lfloor r/8\rfloor+\tfrac12,\qquad
 B=\lfloor4r\rfloor+\tfrac12,
\]
and the exact finite-window Fourier integrals
\begin{equation}\label{pcn:fs:eq:sectors}
 I_j(n)=\int_a^B Q(t)t^n\ee^{2\pi ijt}\dd t,
 \qquad j\in\mathbb Z.
\end{equation}
For every fixed $j\ge0$, let
\begin{equation}\label{pcn:fs:eq:saddles}
 w_j=W_{-j}(n/\ee),\qquad t_j=\frac n{w_j},\qquad
 M_j=\frac{\ee^{-2}t_j^{n-3/2}\exp(t_j-n)}
                 {\sqrt{2\pi(w_j+1)}},\qquad E_j=|M_j|.
\end{equation}
The square root is the continuation of the positive root, equivalently the principal one for sufficiently large $n$. In particular $M_0=E_0$ is the real leading equivalent proved in Part~\ref{pcn:part:core}: it is $M_n$ of \eqref{pcn:eq:rw}, and Theorem~\ref{pcn:thm:main} gives $a_n\sim M_n$.

\begin{theorem}\label{pcn:fs:thm:main}
For every fixed integer $J\ge0$ and every $0<d<\log2$,
\begin{equation}\label{pcn:fs:eq:decomp}
 a_n=I_0(n)+2\Re\sum_{j=1}^J I_j(n)
       +O_J(E_{J+1})+O_{J,d}(M_0\ee^{-dn}).
\end{equation}
For every fixed $j\ge0$ and $L\ge1$,
\begin{equation}\label{pcn:fs:eq:sector-expansion}
 I_j(n)=M_j\left\{\sum_{\ell=0}^{L-1}
                 \frac{c_\ell(w_j)}{t_j^\ell}
                    +O_{j,L}(r^{-L})\right\},
\end{equation}
where the rational functions $c_\ell$ are exactly those of the real saddle in Part~\ref{pcn:part:core} (Theorem~\ref{pcn:thm:main} and the algorithm \eqref{pcn:eq:c-algorithm}). The first two are
\[
 c_0(u)=1,\qquad
 c_1(u)=\frac{20u^3+86u^2+128u+67}{24(u+1)^3}.
\]
For every fixed $j\ge1$,
\begin{equation}\label{pcn:fs:eq:envelope}
 \log\frac{E_j}{M_0}
 =-\frac{2\pi^2j^2r}{w+1}+O_j(r/w^3)
 =-\frac{2\pi^2j^2n}{w^2}\bigl(1+O_j(w^{-1})\bigr).
\end{equation}
Thus $E_{j+1}=o(E_j)$, and $M_0\ee^{-dn}=o(E_j)$ for every fixed $d>0$.
\end{theorem}

The proof separates three errors. Spectral replacement costs $O_d(M_0\ee^{-dn})$ for all $d<\log2$. The real-window and paired vertical-edge errors are $O(M_0\ee^{-\gamma n})$ with a fixed $\gamma>\log2$, for example $\gamma=1$. The remaining horizontal Fourier sum is $O_J(E_{J+1})$. Constants can depend on the fixed mode cutoff and truncation order.

\begin{remark}\label{pcn:fs:rem:scope}
The real pair $2\Re I_j$ is oscillatory and can vanish; its error is measured against the positive envelope $E_j$, not relative to the real pair. Formula \eqref{pcn:fs:eq:decomp} involves the exact Bessel integrals \eqref{pcn:fs:eq:sectors}. A fixed algebraic truncation of $I_0$ has error $O(M_0r^{-L})$, larger than every fixed nonzero Fourier envelope. Therefore substituting finitely many terms into all retained sectors does not by itself give an exponentially accurate numerical approximation.
\end{remark}

The statements also hold for real $n\to\infty$ if $a_n$ means the canonical positive spectral interpolation \eqref{pcn:eq:interpolation} of Part~\ref{pcn:part:core}. This extension is a statement about counts; no exponentially accurate inverse is claimed here.

\section{Global exponentially accurate spectral replacement}
The audited spectrum of Part~\ref{pcn:part:core} gives (Proposition~\ref{pcn:prop:nodes}, equations \eqref{pcn:eq:nodeerror}, \eqref{pcn:eq:weight-b}, \eqref{pcn:eq:weight-leading} and \eqref{pcn:eq:spectrum}), for every fixed $0<c<1$ and all sufficiently large integers $k$,
\begin{align}
 |\lambda_k-k|&\le\ee^{-ck\log k},\qquad
 \left|\frac{\omega_k}{Q(k)}-1\right|
     \le C_c\ee^{-ck\log k},\label{pcn:fs:eq:spectral-input}\\
 0<Q(k)&\le Ck^{-2}\exp(-k\log k+2k),\qquad
 a_n=\sum_\lambda\omega_\lambda\lambda^n.\label{pcn:fs:eq:weight-bound}
\end{align}
All weights are positive, their sum is one, and the finitely many excluded nodes lie in a bounded interval. Fixed multiplicative constants in the first estimate may be absorbed by decreasing $c$.

\begin{proposition}\label{pcn:fs:prop:replace}
For a sufficiently large fixed integer $K$, and every $0<d<\log2$,
\[
 a_n=\sum_{k\ge K}Q(k)k^n+O_d(M_0\ee^{-dn}).
\]
\end{proposition}
\begin{proof}
Split at $k=\sqrt n$. By positivity and total mass one, all lower spectral nodes together contribute at most $(2\sqrt n)^n$ for large $n$. The corresponding lower lattice sum has the same bound up to a fixed constant, because $\sum_{k\ge K}Q(k)<\infty$. Since
\[
 \log M_0=n\log n-n\log\log n-n
                +O(n\log\log n/\log n),
\]
these terms are $M_0\exp[-(1/2+o(1))n\log n]$.

For $k\ge\sqrt n$, $n|\lambda_k-k|/k\to0$ uniformly. The mean-value estimate for $(\lambda_k/k)^n$, together with \eqref{pcn:fs:eq:spectral-input}, gives
\[
 |\omega_k\lambda_k^n-Q(k)k^n|
 \le C_c(n+1)Q(k)k^n\ee^{-ck\log k}.
\]
Put $p=1+c$ and $F_p(x)=n\log x-px\log x+2x$. This is strictly concave and has its unique maximum at
\[
 \frac n s=p\log s+p-2,\qquad
 s=\frac n{pW_0((n/p)\ee^{1-2/p})}.
\]
Consequently $s/r=p^{-1}(1+O_p(w^{-1}))$ and
\begin{equation}\label{pcn:fs:eq:modified-saddle}
 F_p(s)=n\log s-n+ps,\qquad
 F_p(s)-F_1(r)=-n\log p+O_p(n/w).
\end{equation}
Up to $n^2$ there are at most $n^2$ summands, each bounded by $\exp F_p(s)$. For $k>n^2$, eventually $F_p(k)\le-(p/2)k\log k$, so that tail is smaller than one. Polynomial factors and the prefactor between $\exp F_1(r)$ and $M_0$ are harmless. Thus the full error is
\[
 M_0O_c\left(\exp[-n\log(1+c)+O_c(n/w)]\right).
\]
For a given $d<\log2$, choose $c<1$ such that $\log(1+c)>d$. This proves the claim globally, rather than only near the principal saddle.
\end{proof}

\section{The complex Bessel weight}
\begin{lemma}\label{pcn:fs:lem:wedge}
For every fixed $0<\theta<\pi/2$ and $N\ge1$, uniformly as $|z|\to\infty$ in $|\arg z|\le\theta$,
\begin{equation}\label{pcn:fs:eq:b-wedge}
 b(z)=-\frac\ee\pi z\Gamma(z)z^{-z/2}
       \left(\sum_{m=0}^{N-1}\mathcal B_mz^{-m}
                    +O_{N,\theta}(|z|^{-N})\right),
\end{equation}
where $\mathcal B_0=1$. In particular $b$ is eventually zero-free there, and
\begin{equation}\label{pcn:fs:eq:Q-wedge}
 Q(z)=\frac{z^{-2}}{2\pi\ee^2}\ee^{-z\Log z+2z}
        \left(\sum_{m=0}^{N-1}\alpha_mz^{-m}
                     +O_{N,\theta}(|z|^{-N})\right).
\end{equation}
The coefficients agree with the real expansions, including
\[
 \mathcal B_1=-\tfrac12,\quad\mathcal B_2=-\tfrac5{24},\qquad
 \alpha_0=1,\quad\alpha_1=\tfrac56,\quad\alpha_2=\tfrac{73}{72}.
\]
Each order-$N$ normalized remainder is holomorphic; on a smaller closed wedge its $q$th derivative is $O_{N,q,\theta}(|z|^{-N-q})$.
\end{lemma}

\subsection{Schl\"afli representation and contour rotation}
Write $z=R\ee^{i\phi}=u+iv$, so $u\ge R\cos\theta$. Schl\"afli's exact formula \cite[10.9.7]{DLMF} is
\begin{align*}
 Y_\nu(\zeta)={}&\frac1\pi\int_0^\pi
       \sin(\zeta\sin s-\nu s)\dd s\\
 &-\frac1\pi\int_0^\infty
       \left(\ee^{\nu s}+\cos(\pi\nu)\ee^{-\nu s}\right)
                                      \ee^{-\zeta\sinh s}\dd s,
 \qquad |\arg\zeta|<\pi/2.
\end{align*}
All parameter integrals here converge locally uniformly in $\Re z>0$. Combining the growing terms and putting $\xi=\sqrt z\,\ee^s$ gives
\begin{equation}\label{pcn:fs:eq:schlafli}
 b(z)=-\frac{z^{-z/2}}\pi\int_{\sqrt z}^{\infty\ee^{i\phi/2}}
         (\xi-1)\xi^{z-1}\ee^{-\xi+z/\xi}\dd\xi+E_{\mathrm S}(z),
\end{equation}
where
\[
 |E_{\mathrm S}(z)|\le C_\theta(1+\sqrt R)\ee^{\pi|v|+2\sqrt R}.
\]
Indeed the trigonometric integral has this bound, and the decaying integrals are bounded by $1/u$ or $1/(u+1)$ before their explicit factors.
Uniform Stirling gives, for $G(z)=z\Gamma(z)z^{-z/2}$,
\begin{equation}\label{pcn:fs:eq:G-size}
 |G(z)|\asymp_\theta R^{1/2}
   \exp\left(\tfrac12u\log R-\tfrac12v\phi-u\right).
\end{equation}
Thus $E_{\mathrm S}/G=O(\exp[-\eta R\log R])$ for any fixed
$\eta<\tfrac12\cos\theta$.

Rotate the integration ray in \eqref{pcn:fs:eq:schlafli} from argument $\phi/2$ to argument $\phi$, retaining starting radius $\sqrt R$. The large closing arc tends to zero since $\Re\xi\ge|\xi|\cos\theta>0$. No cut or singularity is crossed. On the small connecting arc $\xi=\sqrt R\ee^{i\psi}$,
\[
 |z^{-z/2}\xi^{z-1}|=R^{-1/2}\ee^{v(\phi/2-\psi)}
                         \le R^{-1/2}.
\]
Including $\xi-1$, the arc length, and $|\ee^{-\xi+z/\xi}|\le\ee^{2\sqrt R}$ bounds the arc by $C_\theta\sqrt R\ee^{2\sqrt R}$ after the $z^{-z/2}$ normalization. It is again superexponentially small relative to \eqref{pcn:fs:eq:G-size}.

The rotated integral, denoted $\mathcal I(z)$, becomes
\begin{equation}\label{pcn:fs:eq:rotated}
 \mathcal I(z)=z^z\int_{R^{-1/2}}^\infty
                 (zx-1)x^{z-1}\ee^{-zx+1/x}\dd x.
\end{equation}
This rotation is important: on the positive $x$ ray,
$|x^z\ee^{-zx}|=\exp[u(\log x-x)]$, so the absolute-value saddle is exactly at $x=1$.

\subsection{Uniform Gamma normalization and tails}
Let $h(x)=x-1-\log x$. Stirling also gives
\[
 |z\Gamma(z)|\asymp_\theta R^{1/2}|z^z|\ee^{-u}.
\]
Since $h(x)\ge c(x-1)^2$ on $[1/2,2]$,
\begin{equation}\label{pcn:fs:eq:central-norm}
 \frac{|z^z|}{|z\Gamma(z)|}
       \int_{1/2}^2|zx-1|x^{u-1}\ee^{-ux}\dd x\le C_\theta.
\end{equation}
On $R^{-1/2}\le x\le1/2$, use $\ee^{1/x}\le\ee^{\sqrt R}$ and
\[
 \int_0^{1/2}x^{u-1}\ee^{-ux}\dd x
       \le\frac2u\ee^{-u}\ee^{-u h(1/2)}.
\]
This follows by the substitution $h$, with
$\dd x/x=-\dd h/(1-x)$ and $(1-x)^{-1}\le2$. After normalization, the lower tail is $O_\theta(\ee^{-c_\theta R})$.
For $x\ge2$, use $\ee^{1/x}\le\ee^{1/2}$, $|zx-1|/x\le R+1$, and $h'(x)\ge1/2$. The upper tail has the same exponential bound. Hence \eqref{pcn:fs:eq:rotated}, divided by $z\Gamma(z)$, can be restricted to $[1/2,2]$ with exponential error.

\subsection{Finite Gamma sums and a summable product remainder}
Take $H=\lfloor R^{1/4}\rfloor$. On $[1/2,2]$,
\[
 \ee^{1/x}=\sum_{j=0}^H\frac{x^{-j}}{j!}
             +O\left(\frac{\ee^2 2^{H+1}}{(H+1)!}\right).
\]
By \eqref{pcn:fs:eq:central-norm}, the normalized integral error is smaller than every fixed power of $R^{-1}$. Extend each retained polynomial term to $(0,\infty)$; do not extend the singular untruncated $\ee^{1/x}$. For $j\le H<u/4$,
\[
 \int_0^{1/2}x^{u-j-1}\ee^{-ux}\dd x
 \le\frac{2^j}{u/2-j}\ee^{-u}\ee^{-u h(1/2)}.
\]
To verify this, differentiate $x^{u-j}\ee^{-ux}$ and use
$u-j-ux\ge u/2-j$. Summing against $1/j!$ costs at most $\ee^2$. On the upper tail $x^{-j}\le1$, so summation costs at most $\ee$. All extension errors are exponentially small.

The complex Laplace--Gamma identity now gives
\[
 z^z\int_0^\infty(zx-1)x^{z-j-1}\ee^{-zx}\dd x
       =z^j\{\Gamma(z+1-j)-\Gamma(z-j)\}.
\]
It is valid since $\Re(z-j)>0$ and $\Re z>0$, with the compatible principal powers. Thus, for every fixed $M$,
\begin{align}\label{pcn:fs:eq:product-sum}
 \frac{\mathcal I(z)}{z\Gamma(z)}
 =\sum_{j=0}^H\frac1{j!}\left\{
 \prod_{\ell=0}^{j-1}(1-\ell/z)^{-1}
 -\frac1z\prod_{\ell=1}^{j}(1-\ell/z)^{-1}\right\}
 +O_{M,\theta}(R^{-M}).
\end{align}
Empty products equal one. For $\delta=0,1$, put
\[
 P_{j,\delta}(s)=\prod_{\ell=\delta}^{j-1+\delta}(1-\ell s)^{-1}
                =\sum_{m\ge0}p_{m,\delta}(j)s^m.
\]
For fixed $m$, $p_{m,\delta}(j)$ is a polynomial in $j$ of degree at most $2m$, by expanding the logarithm and using power sums. On
$|s|=\rho_j=[4(j+1)^2]^{-1}$, $|P_{j,\delta}(s)|\le\ee$. For $j\le H$, eventually $R^{-1}\le\rho_j/2$. Cauchy's coefficient bound and a geometric sum give
\begin{equation}\label{pcn:fs:eq:product-remainder}
 \left|P_{j,\delta}(1/z)-\sum_{m=0}^{N-1}
                  p_{m,\delta}(j)z^{-m}\right|
       \le2\ee\,4^N(j+1)^{2N}R^{-N}.
\end{equation}
This is summable against $1/j!$, uniformly in $R$. Extending each coefficient sum from $j\le H$ to all $j$ costs another factorial tail. We obtain
\[
 \frac{\mathcal I(z)}{z\Gamma(z)}
       =\ee\sum_{m=0}^{N-1}\mathcal B_mz^{-m}+O_{N,\theta}(R^{-N}),
\]
where the coefficientwise generator is
\begin{equation}\label{pcn:fs:eq:beta-generator}
 \mathcal B(s)=\ee^{-1}\sum_{j\ge0}\frac1{j!}
 \left\{\prod_{\ell=0}^{j-1}(1-\ell s)^{-1}
      -s\prod_{\ell=1}^{j}(1-\ell s)^{-1}\right\}
        =\sum_{m\ge0}\mathcal B_ms^m.
\end{equation}
Equation \eqref{pcn:fs:eq:beta-generator} means coefficientwise summation, not convergence as an analytic function of $s$.

\subsection{Zero-freeness and analytic remainder bounds}
Combining the preceding estimates proves \eqref{pcn:fs:eq:b-wedge}. With $N=1$, its normalized factor is $1+O(1/R)$, so its modulus is eventually at least $1/2$. The remaining factors do not vanish. This proves eventual zero-freeness.
Algebraically,
\[
 Q(z)=\ee^{-2}\frac{z^{z-3}}{\Gamma(z)^2}
          \left(\sum_{m=0}^{N-1}\mathcal B_mz^{-m}+O(R^{-N})\right)^{-2}.
\]
Uniform Stirling \cite[\S5.11]{DLMF} proves \eqref{pcn:fs:eq:Q-wedge}, with formal coefficient generator
\begin{equation}\label{pcn:fs:eq:alpha-generator}
 \sum_{m\ge0}\alpha_ms^m=
 \mathcal B(s)^{-2}\exp\left[-\sum_{q\ge1}
                \frac{\beta_{2q}s^{2q-1}}{q(2q-1)}\right].
\end{equation}
For derivative estimates, choose a larger wedge with angle still below $\pi/2$. A disk of radius $\epsilon|z|$ centered in the smaller wedge lies in the larger one, with all moduli comparable to $|z|$. Cauchy's formula gives the asserted derivative bounds for the analytic normalized remainder. The nonanalytic cutoff $H$ was used only to establish the bound and is never differentiated. In particular,
\begin{equation}\label{pcn:fs:eq:log-derivative}
 \frac{Q'(z)}{Q(z)}=1-\Log z-\frac2z+O(|z|^{-2}).
\end{equation}
This completes Lemma~\ref{pcn:fs:lem:wedge}.

\section{Poisson summation and the entire omitted Fourier tail}
\subsection{The real window}
Let $\phi(x)=n\log x-x\log x+2x$. For fixed $c>0$,
\begin{equation}\label{pcn:fs:eq:window-rate}
 \phi(cr)-\phi(r)=-n\mathcal H(c)
                +r\{c(1-\log c)-1\},\qquad
 \mathcal H(c)=c-1-\log c.
\end{equation}
Here $\mathcal H(1/8)=\log8-7/8>1.20$ and
$\mathcal H(4)=3-\log4>1.61$. Outside the window, $\phi'$ has the appropriate sign and magnitude at least a positive constant times $w$. Concavity bounds each discrete tail by a geometric sum from its boundary. Polynomial amplitudes and $O(1)$ endpoint rounding cost only $\exp(o(n))$. Therefore, for any fixed
$\gamma<\min(\mathcal H(1/8),\mathcal H(4))$,
\begin{equation}\label{pcn:fs:eq:window-tail}
 \sum_{k\ge K}Q(k)k^n
       =\sum_{a<k<B}Q(k)k^n+O(M_0\ee^{-\gamma n}).
\end{equation}
We may and will choose $\gamma=1$.

Put $f(z)=Q(z)z^n$. Finite-interval Poisson summation gives
\begin{equation}\label{pcn:fs:eq:poisson}
 \sum_{a<k<B}f(k)=I_0+2\sum_{j\ge1}\Re I_j.
\end{equation}
The paired series converges absolutely. Indeed, the first integration by parts produces the boundary term
\[
 \frac{(-1)^j}{2\pi i j}\{f(B)-f(a)\},
\]
which is purely imaginary. A second integration by parts bounds the remaining real part by $C_n/j^2$. Alternatively, ordinary Fourier convergence of the piecewise smooth periodicization at zero proves \eqref{pcn:fs:eq:poisson}; no jump occurs there since both endpoints are half-integers.

\subsection{One common contour height}
Fix $m=J+1\ge1$, and write
\[
 \psi_j(z)=n\Log z-z\Log z+2z+2\pi ijz.
\]
The saddle equation is
\[
 n/t_j=\Log t_j-1-2\pi ij=w_j,
\]
with branch specified by $w_j+\Log w_j=\log(n/\ee)-2\pi ij$. For fixed $j$,
\[
 w_j=w-2\pi ij+O_j(w^{-1}),\qquad
 t_j/r=1+O_j(w^{-1}).
\]
For $j\ge1$, $\Im t_j>0$. In particular
\[
 h=\Im t_m\sim\frac{2\pi m r}{w+1}\longrightarrow\infty.
\]
Every rectangle with real part between $a$ and $B$ and height $h$ lies eventually in any fixed positive-angle wedge. Lemma~\ref{pcn:fs:lem:wedge} therefore permits shifting every mode $j\ge m$ to this same height. Cauchy's theorem, with its orientation retained, gives
\begin{align}
 I_j={}&H_j+i(-1)^j\int_0^h\ee^{-2\pi jy}
                    \{f(a+iy)-f(B+iy)\}\dd y,\label{pcn:fs:eq:rectangle}\\
 H_j={}&\int_a^B f(x+ih)\ee^{2\pi ij(x+ih)}\dd x.\nonumber
\end{align}

\subsection{Paired vertical endpoints}
Since $f(a)$ and $f(B)$ are real, for either endpoint $x$,
\[
 |\Im f(x+iy)|\le y\sup_{0\le v\le h}|f'(x+iv)|.
\]
For $0\le v\le h$ and $x$ a fixed positive multiple of $r$,
\[
 \Re\phi(x+iv)-\phi(x)=O_J(n/w^2).
\]
This follows directly from $v/x=O_J(w^{-1})$ in the analytic phase. Equations \eqref{pcn:fs:eq:Q-wedge} and \eqref{pcn:fs:eq:log-derivative} give $|f'|\le C_Jw|f|$ on the edges. Combining with \eqref{pcn:fs:eq:window-rate}, and absorbing subexponential factors, gives
\[
 \sup_{0\le v\le h}|f'(x+iv)|\le M_0\ee^{-\gamma n}
\]
for every fixed $\gamma<\min(\mathcal H(1/8),\mathcal H(4))$ and sufficiently large $n$.
The sum of the absolute real parts of the vertical pieces in \eqref{pcn:fs:eq:rectangle} is finite, because
\begin{equation}\label{pcn:fs:eq:kernel}
 \int_0^h\frac{y\ee^{-2\pi my}}{1-\ee^{-2\pi y}}\dd y
 \le\int_0^\infty\frac{y\ee^{-2\pi my}}{1-\ee^{-2\pi y}}\dd y
 <\infty.
\end{equation}
Thus the complete paired vertical tail is $O_J(M_0\ee^{-\gamma n})$. The factor $y$ from real pairing removes the $1/y$ singularity at zero. No absolutely convergent unsigned series of the original complex $I_j$ is asserted.

\subsection{A global horizontal Gaussian bound}
The horizontal terms do admit an absolute geometric bound:
\begin{equation}\label{pcn:fs:eq:horizontal-sum}
 \sum_{j\ge m}|H_j|
 \le\frac1{1-\ee^{-2\pi h}}
       \int_a^B|f(x+ih)|\ee^{-2\pi mh}\dd x.
\end{equation}
At $x_m=\Re t_m$, the derivative of $\Re\psi_m(x+ih)$ vanishes. Throughout the full window,
\begin{equation}\label{pcn:fs:eq:curvature}
 \Re\psi_m''(x+ih)
 =-\frac{n(x^2-h^2)}{(x^2+h^2)^2}
       -\frac{x}{x^2+h^2}\le-c_Jw/r
\end{equation}
for all sufficiently large $n$, since $h/a\to0$. The saddle is therefore the unique global maximum on this segment. Since $|x+ih|\asymp r$, Lemma~\ref{pcn:fs:lem:wedge} and the Gaussian bound yield
\[
 \int_a^B|f(x+ih)|\ee^{-2\pi mh}\dd x
 \le C_Jr^{-2}\ee^{\Re\psi_m(t_m)}\sqrt{r/w}
 \le C_JE_m.
\]
The prefactor in \eqref{pcn:fs:eq:horizontal-sum} tends to one. Proposition~\ref{pcn:fs:prop:replace}, \eqref{pcn:fs:eq:window-tail}, and \eqref{pcn:fs:eq:poisson}--\eqref{pcn:fs:eq:curvature} prove \eqref{pcn:fs:eq:decomp}. This controls infinitely many omitted modes without summing infinitely many saddle expansions.

\section{The all orders expansion of each exact sector}
For fixed $j\ge1$, shift only $I_j$ to height $\Im t_j$. Its two vertical pieces separately are $O(M_0\ee^{-\gamma n})$; the preceding endpoint bound suffices without pairing. For $j=0$, stay on the real axis. On either horizontal segment, the same global Gaussian localization applies.

Near the saddle, write $z=t_j+u$ with $u$ real. Then
\[
 A_j=-\psi_j''(t_j)=\frac{w_j+1}{t_j},\qquad
 \Re A_j\asymp w/r>0,
\]
and for $q\ge3$,
\begin{equation}\label{pcn:fs:eq:derivative-formula}
 \psi_j^{(q)}(t_j)=(-1)^{q-1}(q-2)!
                  \frac{(q-1)w_j+1}{t_j^{q-1}}.
\end{equation}
Here is a uniform fixed-order remainder argument. Put
$\sigma=\sqrt{r/w}$, $\rho=(rw)^{-1/2}$, and $u=\sigma v$.
The quadratic coefficient $A_j\sigma^2$ stays in a compact subset of $\{\Re z>0\}$. The coefficient of $v^{h+2}$ in the nonquadratic exponent is $O_h(\rho^h)$, and $u/t_j=O(\rho v)$.

On $|v|\le\rho^{-\epsilon}$, for sufficiently small fixed $\epsilon>0$, the nonquadratic exponent has size at most $C\rho|v|^3=o(1)$. Expand the phase to sufficiently high degree and expand its exponential by total powers of $\rho$. Retaining weights below $2N$ has remainder bounded by
\[
 C_N\rho^{2N}P_N(|v|)\exp(C\rho|v|^3),
\]
for a fixed polynomial $P_N$. The phase Taylor remainder has the same type of bound by analyticity on $|u|\le\eta r$ and \eqref{pcn:fs:eq:derivative-formula}. These bounds are integrable against $\ee^{-cv^2}$. Outside the central interval, global Gaussian localization gives $\exp(-c\rho^{-2\epsilon})$, smaller than every fixed power of $r^{-1}$ even after polynomial factors.

Apply this to the amplitude term
$\alpha_\ell t_j^{-\ell}(1+u/t_j)^{-\ell-2}$ with $N=L-\ell$.
Its error is
\[
 O\bigl(r^{-\ell}\rho^{2(L-\ell)}\bigr)
       =O\bigl(r^{-L}w^{-(L-\ell)}\bigr)=O(r^{-L}).
\]
The remainder of \eqref{pcn:fs:eq:Q-wedge} itself contributes $O(r^{-L})$ by the absolute Gaussian bound. Odd total weights have odd degree in $v$ and vanish on symmetric extension to the full real line; that extension has superalgebraically small error. Summing over the finitely many $\ell<L$ establishes the relative remainder in \eqref{pcn:fs:eq:sector-expansion}.

For $\Re A>0$, analytic continuation of real Gaussian moments gives
\[
 \int_{\mathbb R}u^{2q}\ee^{-Au^2/2}\dd u
       =\sqrt{2\pi}(2q-1)!!A^{-q-1/2},
\]
with zero odd moments. Substitution reproduces precisely the real coefficient algorithm evaluated at $w_j$. The leading factor is
\[
 \frac{t_j^{-2}}{2\pi\ee^2}\ee^{\psi_j(t_j)}
       \sqrt{\frac{2\pi t_j}{w_j+1}}=M_j,
\]
since $\psi_j(t_j)=n\Log t_j-n+t_j$. For each fixed $j,L$, the endpoint errors are smaller than $E_jr^{-L}$ and are absorbed into the same remainder.

For completeness, use formal Gaussian expectation with
$\E V^{2q}=(2q-1)!!/(u+1)^q$ and $\E V^{2q+1}=0$. Then
\begin{align}\label{pcn:fs:eq:coefficient-generator}
 c_\ell(u)=[s^{2\ell}]\E\Bigg[&
 \sum_{k\ge0}\alpha_ks^{2k}(1+Vs)^{-k-2}\\
 &\times\exp\left\{\sum_{h\ge1}
 \frac{(-1)^{h+1}((h+1)u+1)V^{h+2}s^h}{(h+2)(h+1)}\right\}\Bigg].\nonumber
\end{align}
Only finitely many terms enter each coefficient. In particular,
\[
 c_2(u)=\frac{1168u^6+9968u^5+34692u^4+64536u^3
              +69380u^2+42208u+11857}{1152(u+1)^6}.
\]
The supplied symbolic generator computes these coefficients; there is no numerical fitting.

\section{Ordering the envelopes}
For bounded real $v$, let $t(v)$ solve
\[
 n/t=\Log t-1-iv
\]
by continuation from $t(0)=r$. Set $S(v)=n\Log t(v)-n+t(v)$, the stationary value of the phase including $ivt$. Differentiation gives
\[
 S'(v)=it(v),\qquad
 t'(v)=\frac{it(v)}{n/t(v)+1},\qquad
 S''(0)=-\frac r{w+1}.
\]
Repeated differentiation of the saddle relation yields
$S^{(q)}(v)=O_q(r/w^{q-1})$ for $q\ge2$, uniformly on bounded $v$ intervals. Conjugation makes $\Re S$ even. Taylor's theorem through degree three therefore gives
\[
 \Re\{S(2\pi j)-S(0)\}
       =-\frac{2\pi^2j^2r}{w+1}+O_j(r/w^3).
\]
The logarithm of the absolute prefactor ratio in \eqref{pcn:fs:eq:saddles} is $O_j(w^{-2})$, absorbed into the displayed remainder. This proves \eqref{pcn:fs:eq:envelope}, completes Theorem~\ref{pcn:fs:thm:main}, and shows that the spectral replacement error lies below every fixed Fourier sector.

\section{Diagnostics and scope}\label{pcn:fs:sec:scope}
The reproducibility package contains two numerical checks. The first evaluates the exact Bessel weight on central locally steepest saddle segments and compares the result with the first three rational corrections. With $160$-node Gauss--Legendre quadrature at $70$ decimal digits, the difference between the third scaled remainder and $c_3(w_j)$ is approximately
\begin{center}
\begin{tabular}{rrr}
\toprule
$n$ & $j$ & Scaled remainder discrepancy\\
\midrule
$10^4$ & $0$ & $4.89\times10^{-3}$\\
$10^4$ & $1$ & $5.30\times10^{-3}$\\
$3\times10^4$ & $0$ & $1.76\times10^{-3}$\\
$3\times10^4$ & $1$ & $1.90\times10^{-3}$\\
$10^5$ & $0$ & $5.76\times10^{-4}$\\
$10^5$ & $1$ & $6.17\times10^{-4}$\\
$10^8$ & $2$ & $9.98\times10^{-7}$\\
\bottomrule
\end{tabular}
\end{center}
These non-certified diagnostics integrate central saddle segments; they do not evaluate complete cutoff-defined sectors. The proof controls the omitted contour pieces asymptotically. Separate checks recover $\alpha_3=11821/6480$ along arguments $0$, $0.4$, and $1.2$ radians, and compare exact recurrence counts with a truncated exact-Bessel lattice sum.

Poisson summation, complex Gaussian expansion, and Lambert-$W$ saddle descriptions in \mbox{Bell/Touchard} asymptotics are standard methods; see Paris~\cite{Paris} for complex Touchard saddle expansions. The work here verifies the necessary estimates for the exact powered-Catalan Bessel weight, including its zero-free wedge and the entire omitted Fourier tail. No comprehensive claim of novelty or priority is made.
\begin{writenote}
The repository's transseries volume, \repopath{Analysis/Transseries/docs/series-and-transseries/Transseries_And_Inversion/}, treats the Bell numbers ($\theta=1$) by the same Lambert saddle (theorem \texttt{q2:thm:bell}) and, in remark \texttt{q2:rem:bell-sectors}, keeps their nonprincipal complex saddles \emph{formal}, because identifying the Stokes multipliers of the Cauchy contour is a separate global problem. This Part proves fixed complex sectors rigorously for a different sequence, A113227, starting from an exact positive spectrum and a lattice sum rather than from a Cauchy contour. It does not settle the Bell-number question of that remark, and nothing here is claimed as new about the method.
\end{writenote}

The result fixes both the number of sectors and the algebraic order. It neither sums infinitely many asymptotic sector expansions nor proves a theorem with $J$ growing with $n$. The sectors depend on the explicit cutoff convention; other sufficiently wide fixed relative windows change them by an exponentially smaller endpoint contribution. No canonical Borel sum, convergence of the infinite correction series, or exponential inverse theorem follows from the arguments above. The algebraic inverse of Part~\ref{pcn:part:core} (Theorem~\ref{pcn:thm:inverse}) remains valid independently.

\paragraph{Review status.}
The spectrum and real all-orders input was independently audited. Both new proof components received independent cross-review, followed by integrated mathematical review. These checks are not formal verification or conventional peer review. The computations supplement the proof and do not replace it.
\begin{writenote}
The three programs and their outputs are shipped as \repopath{code/56-fourier-*} and \repopath{data/56-fourier-*}; the coefficient generators that source 56's replay also calls are Part~\ref{pcn:part:core}'s (\repopath{code/57-catalan-residue_expansion.py} and \repopath{code/57-catalan-saddle_expansion.py}), of which source 56 carried byte-identical copies. At intake (2 October 2026) \texttt{verify\_fourier\_sectors.py} and \texttt{verify\_exponential\_spectral\_replacement.py} reproduced their delivered outputs; the order-192 refinement and the consistency check were not rerun (README).
\end{writenote}


\begingroup\small
\begin{thebibliography}{99}
\setlength{\itemsep}{3pt}
\bibitem{OEIS} OEIS Foundation, \emph{A113227}, entry inspected 1 October 2026. Formal continued fraction credited to S.\ N.\ Gladkovskii, 10 October 2012. \url{https://oeis.org/A113227}.
\bibitem{Elizalde} S.\ Elizalde, \emph{Asymptotic enumeration of permutations avoiding generalized patterns}, Advances in Applied Mathematics \textbf{36} (2006), 138--155. \href{https://doi.org/10.1016/j.aam.2005.05.006}{doi:10.1016/j.aam.2005.05.006}; \href{https://arxiv.org/abs/math/0505254}{arXiv:math/0505254}.
\bibitem{Callan} D.\ Callan, \emph{A bijection to count $(1$-$23$-$4)$-avoiding permutations} (2010). \href{https://arxiv.org/abs/1008.2375}{arXiv:1008.2375}.
\bibitem{Beaton} N.\ R.\ Beaton, M.\ Bouvel, V.\ Guerrini and S.\ Rinaldi, \emph{Enumerating five families of pattern-avoiding inversion sequences; and introducing the powered Catalan numbers}, Theoretical Computer Science \textbf{777} (2019), 69--92. \href{https://doi.org/10.1016/j.tcs.2019.02.003}{doi:10.1016/j.tcs.2019.02.003}; \href{https://arxiv.org/abs/1808.04114}{arXiv:1808.04114}.
\bibitem{Duchi} E.\ Duchi, V.\ Guerrini and S.\ Rinaldi, \emph{A generating tree for permutations avoiding the pattern $122^+3$}, Fundamenta Informaticae \textbf{163} (2018), 21--39. \href{https://doi.org/10.3233/FI-2018-1730}{doi:10.3233/FI-2018-1730}.
\bibitem{LinFu} Z.\ Lin and S.\ Fu, \emph{On $\underline{12}0$-avoiding inversion and ascent sequences}, European Journal of Combinatorics \textbf{93} (2021), 103282. \href{https://doi.org/10.1016/j.ejc.2020.103282}{doi:10.1016/j.ejc.2020.103282}; \href{https://arxiv.org/abs/2003.11813}{arXiv:2003.11813}.
\bibitem{Cervetti} M.\ Cervetti, \emph{A generating tree with a single label for permutations avoiding the vincular pattern $1$-$32$-$4$}, Pure Mathematics and Applications \textbf{30} (2022), 56--62. \href{https://doi.org/10.2478/puma-2022-0009}{doi:10.2478/puma-2022-0009}.
\bibitem{Frosini} A.\ Frosini, V.\ Guerrini and S.\ Rinaldi, \emph{Constrained underdiagonal paths and pattern avoiding permutations}, Mathematics \textbf{13} (2025), 517. \href{https://doi.org/10.3390/math13030517}{doi:10.3390/math13030517}.
\bibitem{Testart} B.\ Testart, \emph{Generating trees growing on the left for pattern-avoiding inversion sequences}, Discrete Mathematics and Theoretical Computer Science \textbf{27:1} (2025), article 10. \href{https://arxiv.org/abs/2411.05726}{arXiv:2411.05726}.
\bibitem{PSZ} M.\ P\'etr\'eolle, A.\ D.\ Sokal and B.-X.\ Zhu, \emph{Lattice paths and branched continued fractions: an infinite sequence of generalizations of the Stieltjes--Rogers and Thron--Rogers polynomials, with coefficientwise Hankel-total positivity}, version 2 (2020), Theorem 9.9. \href{https://arxiv.org/abs/1807.03271}{arXiv:1807.03271}.
\bibitem{DLMF} F.\ W.\ J.\ Olver et al., eds., \emph{NIST Digital Library of Mathematical Functions}. Part~\ref{pcn:part:core}: Chapter 10, especially \S\S10.2, 10.4, 10.5 and formula 10.9.7, \url{https://dlmf.nist.gov/10}. Part~\ref{pcn:part:fs}: \S\S4.13, 5.11, 10.9, in particular Schl\"afli's integral \url{https://dlmf.nist.gov/10.9.E7}, complex Stirling estimates \url{https://dlmf.nist.gov/5.11}, and Lambert branches \url{https://dlmf.nist.gov/4.13}. [Merged from the two sources' entries.]
\bibitem{Paris} R.\ B.\ Paris, \emph{The asymptotics of the Touchard polynomials}, 2016, \url{https://arxiv.org/abs/1606.07883}.
\end{thebibliography}
\endgroup
\end{document}
